\SourcePage{26}
\SGAsection{4}{概形范畴中的代数结构}
只要不致引起歧义，我们便容许采用以下滥用语言：用 $T$ 表示 $C_{/S}$ 的\Term{object}{对象} $T\xrightarrow f S$，把 $f$（“$T$ 的\Term{structure-morphism}{结构态射}”）的数据略去不写；并把 $C$ 等同于 $\widehat C$ 的一个\Term{subcategory}{子范畴}。根据前几节所述的相容性，这些等同不会带来危险。
我们还按下例简化名称：$(\mathrm{Sch})$-\Term{group}{群}也称为\Term{group-scheme}{群概形}；$(\mathrm{Sch})_{/S}$-\Term{group}{群}称为 $S$ 上的\Term{group-scheme}{群概形}，或 $S$-\Term{group-scheme}{群概形}，或 $S$-\Term{group}{群}；当 $S$ 是\Term{ring}{环} $A$ 的\Term{spectrum}{谱}时，也称为 $A$-\Term{group}{群}。

\textbf{4.1. 常\Term{scheme}{概形}。} \Term{scheme}{概形}\Term{category}{范畴}满足 1.8 所要求的条件。因此可以在其中定义\Term{constant-object}{常对象}：对每个\Term{set}{集合} $E$，都有一个\Term{scheme}{概形} $E_{\mathbb Z}$；对每个\Term{scheme}{概形} $S$，都有一个 $S$-\Term{scheme}{概形} $E_S=(E_{\mathbb Z})_S$。回忆（见该处），对每个 $S$-\Term{scheme}{概形} $T$，$\operatorname{Hom}_S(T,E_S)$ 就是从\Term{topological-space}{拓扑空间} $T$ 到 $E$ 的\Term{locally-constant}{局部常值}映射所成的\Term{set}{集合}。
\Term{functor}{函子} $E\mapsto E_S$ 与\Term{finite}{有限}\Term{projective-limit}{投射极限}交换。特别地，若 $G$ 是\Term{group}{群}，则 $G_S$ 是 $S$-\Term{group-scheme}{群概形}；若 $O$ 是\Term{ring}{环}，则 $O_S$ 是 $S$-\Term{ring-scheme}{环概形}，等等。

\textbf{4.2. \Term{affine}{仿射} $S$-群。} 回忆关于\Term{affine}{仿射} $S$-\Term{scheme}{概形}的若干事实（EGA II，§1）。若 $S$ 的每个\Term{affine-open}{仿射开集}的\Term{inverse-image}{逆像}都是\Term{affine}{仿射}的，就称 $S$-\Term{scheme}{概形} $T$ 在 $S$ 上\Term{affine}{仿射}。这时 $\mathcal O_S$-\Term{algebra}{代数} $f_*(\mathcal O_T)$ 是\Term{quasi-coherent}{拟凝聚}的，记为 $\mathcal A(T)$（$f$ 表示 $T$ 的\Term{structure-morphism}{结构态射}）。反之，对每个\Term{quasi-coherent}{拟凝聚} $\mathcal O_S$-\Term{algebra}{代数} $\mathcal A$，可以对应一个在 $S$ 上\Term{affine}{仿射}的 $S$-\Term{scheme}{概形}，记为 $\operatorname{Spec}(\mathcal A)$。\Term{functor}{函子} $T\mapsto\mathcal A(T)$ 与 $\mathcal A\mapsto\operatorname{Spec}(\mathcal A)$ 给出在 $S$ 上\Term{affine}{仿射}的 $S$-\Term{scheme}{概形}\Term{category}{范畴}与\Term{quasi-coherent}{拟凝聚} $\mathcal O_S$-\Term{algebra}{代数}\Term{category}{范畴}的\Term{opposite-category}{对偶范畴}之间互为\Term{quasi-inverse}{拟逆}的\Term{equivalence}{等价}。
由此，在\Term{affine}{仿射} $S$-\Term{scheme}{概形} $T$ 上给定\Term{algebraic-structure}{代数结构}，\Term{equivalence}{等价}于在\Term{quasi-coherent}{拟凝聚} $\mathcal O_S$-\Term{algebra}{代数}\Term{category}{范畴}的\Term{opposite-category}{对偶范畴}中，在 $\mathcal A(T)$ 上给定相应的结构。特别地，若 $G$ 是在 $S$ 上\Term{affine}{仿射}的 $S$-\Term{group}{群}，则 $\mathcal A(G)$ 带有\Term{augmentation}{增广} $\mathcal O_S$-\Term{bialgebra}{双代数}结构，也就是说，有两个 $\mathcal O_S$-\Term{algebra}{代数}\Term{morphism}{态射}
\[
\Delta:\mathcal A(G)\longrightarrow\mathcal A(G)\otimes_{\mathcal O_S}\mathcal A(G),\qquad
\varepsilon:\mathcal A(G)\longrightarrow\mathcal O_S,
\]
\SourcePage{27}
分别对应于 $S$-\Term{scheme}{概形}\Term{morphism}{态射}
\[
\pi:G\times G\longrightarrow G,\qquad e_G:S\longrightarrow G.
\]
映射 $\Delta$ 与 $\varepsilon$ 满足下列条件（这些条件表达 $G$ 是一个 $S$-\Term{monoid}{幺半群}）：

（HA 1）$\Delta$ \Term{coassociative}{余结合}，即下图交换：
\[
\begin{tikzcd}[column sep=large]
\mathcal A(G)\arrow[r,"{\Delta}"]\arrow[d,"{\Delta}"']&\mathcal A(G)\otimes_{\mathcal O_S}\mathcal A(G)\arrow[d,"{\mathrm{Id}\otimes\Delta}"]\\
\mathcal A(G)\otimes_{\mathcal O_S}\mathcal A(G)\arrow[r,"{\Delta\otimes\mathrm{Id}}"']&\mathcal A(G)\otimes_{\mathcal O_S}\mathcal A(G)\otimes_{\mathcal O_S}\mathcal A(G).
\end{tikzcd}
\]
（HA 2）$\Delta$ 与 $\varepsilon$ 相容，即以下两个复合都是恒等：
\[
\begin{aligned}
\mathcal A(G)&\xrightarrow\Delta\mathcal A(G)\otimes_{\mathcal O_S}\mathcal A(G)\xrightarrow{\mathrm{Id}\otimes\varepsilon}\mathcal A(G)\otimes_{\mathcal O_S}\mathcal O_S\xrightarrow\sim\mathcal A(G),\\
\mathcal A(G)&\xrightarrow\Delta\mathcal A(G)\otimes_{\mathcal O_S}\mathcal A(G)\xrightarrow{\varepsilon\otimes\mathrm{Id}}\mathcal O_S\otimes_{\mathcal O_S}\mathcal A(G)\xrightarrow\sim\mathcal A(G).
\end{aligned}
\]
\footnote{编者注（29）：当然，\Term{inverse-morphism}{求逆态射} $G\to G$ 诱导一个 $\mathcal O_S$-\Term{algebra}{代数}\Term{morphism}{态射} $\tau:\mathcal A(G)\to\mathcal A(G)$，它与 $\Delta$、$\varepsilon$ 一起使 $\mathcal A(G)$ 成为一个 $\mathcal O_S$-\Term{hopf-algebra}{Hopf 代数}。}
借此机会再指出一次：由 $S$-\Term{group}{群}结构的定义可知，要在一个在 $S$ 上\Term{affine}{仿射}的 $S$-\Term{scheme}{概形} $G$ 上给定这样的结构，无须对 $\mathcal A(G)$ 作任何检验；只须对 $S$ 上的每个 $S'$，在 $G(S')$ 上给定一个关于 $S'$ \Term{functorial}{函子性}的\Term{group}{群}结构即可。此注记作适当变更后同样适用于\Term{morphism}{态射}。

\textbf{4.3. 群 $\mathbb G_a$、$\mathbb G_m$ 与环 $\mathbf O$。}

\textbf{4.3.1.} 设 $\mathbf G_a$ 为从 $(\mathrm{Sch})^\circ$ 到 $(\mathrm{Ens})$ 的\Term{functor}{函子}，由
\[
\mathbf G_a(S)=\Gamma(S,\mathcal O_S)
\]
定义，并赋予由\Term{ring}{环} $\Gamma(S,\mathcal O_S)$ 的\Term{additive-group}{加法群}结构所定义的 $\widehat{(\mathrm{Sch})}$-\Term{group}{群}结构。它可以由一个\Term{affine}{仿射}\Term{scheme}{概形}表示，记为 $\mathbb G_a$，因而这是一个\Term{group-scheme}{群概形}：
\[
\mathbb G_a=\operatorname{Spec}\mathbb Z[T].
\]
事实上，$\mathbf G_a(S)=\operatorname{Hom}(S,\mathbb G_a)=\operatorname{Hom}_{\mathrm{Alg.}}(\mathbb Z[T],\Gamma(S,\mathcal O_S))\simeq\Gamma(S,\mathcal O_S)$。

\SourcePage{28}
因此，对每个\Term{scheme}{概形} $S$，都有一个在 $S$ 上\Term{affine}{仿射}的 $S$-\Term{group}{群}，记为 $\mathbb G_{a,S}$；它对应于 $\mathcal O_S$-\Term{bialgebra}{双代数} $\mathcal O_S[T]$，其\Term{diagonal-map}{对角映射}和\Term{augmentation}{增广}由下式定义：
\[
\Delta(T)=T\otimes1+1\otimes T,\qquad\varepsilon(T)=0.
\]
\textbf{4.3.2.} 设 $\mathbf G_m$ 为从 $(\mathrm{Sch})^\circ$ 到 $(\mathrm{Ens})$ 的\Term{functor}{函子}，由
\[
\mathbf G_m(S)=\Gamma(S,\mathcal O_S)^\times
\]
定义，其中 $\Gamma(S,\mathcal O_S)^\times$ 表示\Term{ring}{环} $\Gamma(S,\mathcal O_S)$ 的\Term{invertible-element}{可逆元}素所成的\Term{multiplicative-group}{乘法群}，并赋予其自然的 $\widehat{(\mathrm{Sch})}$-\Term{group}{群}结构。它可以由一个\Term{affine}{仿射}\Term{scheme}{概形}表示，记为 $\mathbb G_m$：
\[
\mathbb G_m=\operatorname{Spec}\mathbb Z[T,T^{-1}]=\operatorname{Spec}\mathbb Z[\mathbb Z],
\]
这里 $\mathbb Z[\mathbb Z]$ 表示\Term{group}{群} $\mathbb Z$ 在\Term{ring}{环} $\mathbb Z$ 上的\Term{group-algebra}{群代数}。事实上，
\[
\mathbf G_m(S)=\operatorname{Hom}_{\mathrm{Alg.}}(\mathbb Z[T,T^{-1}],\Gamma(S,\mathcal O_S))\simeq\Gamma(S,\mathcal O_S)^\times.
\]
因此，对每个\Term{scheme}{概形} $S$，都有一个在 $S$ 上\Term{affine}{仿射}的 $S$-\Term{group}{群}，记为 $\mathbb G_{m,S}$；它对应于 $\mathcal O_S$-\Term{algebra}{代数} $\mathcal O_S[\mathbb Z]$，其\Term{diagonal-map}{对角映射}和\Term{augmentation}{增广}由下式定义：
\[
\Delta(x)=x\otimes x,\qquad\varepsilon(x)=1\quad(x\in\mathbb Z).
\]
\textbf{4.3.3.} 设 $\mathbf O$ 为带有其 $\widehat{(\mathrm{Sch})}$-\Term{ring}{环}结构的\Term{functor}{函子} $\mathbf G_a$。它由\Term{scheme}{概形} $\operatorname{Spec}\mathbb Z[T]$ 表示；当把后者看作带有\Term{ring-scheme}{环概形}结构时，记为 $\mathbb O$。
因此，对每个\Term{scheme}{概形} $S$，$\mathbb O_S=S\times_{\operatorname{Spec}\mathbb Z}\operatorname{Spec}\mathbb Z[T]=\operatorname{Spec}(\mathcal O_S[T])$ 是一个在 $S$ 上\Term{affine}{仿射}的 $S$-\Term{ring-scheme}{环概形}。（注意：在 EGA II，1.7.13 中，$\mathbb O_S$ 记为 $S[T]$。）

\textbf{4.3.3.1.}\footnote{编者注（30）：此段为增补，将在后文使用（参见 II，1.3）。} 对 $\widehat{(\mathrm{Sch})}$ 的每个\Term{object}{对象} $X$，$\mathbf O(X)=\operatorname{Hom}(X,\mathbf O)$ 带有关于 $X$ \Term{functorial}{函子性}的\Term{ring}{环}结构。特别地，若 $X'$ 是\Term{scheme}{概形}，且给定\Term{morphism}{态射} $x:X'\to X$ 与 $f:X\to\mathbf O$（即 $f\in\mathbf O(X)$），则 $f(x)=f\circ x$ 是 $\mathbf O(X')=\Gamma(X',\mathcal O_{X'})$ 的元素。

\begin{SGAresult}{定义}{}
设 $\pi:M\to X$ 是 $\widehat{(\mathrm{Sch})}$ 的一个\Term{morphism}{态射}，且令 $\mathbf O_X=\mathbf O\times X$。若对每个 $X$-\Term{scheme}{概形} $X'$，在 $\operatorname{Hom}_X(X',M)$ 上给定一个关于 $X'$ \Term{functorial}{函子性}的 $\mathbf O(X')$-\Term{module}{模}结构，就称 $M$ 为 $\mathbf O_X$-\Term{module}{模}。
\end{SGAresult}
这\Term{equivalence}{等价}于给定一个 $X$-\Term{abelian-group}{交换群}运算 $\mu:M\times_X M\to M$ 和一个“外运算”
\[
\mathbf O\times M=\mathbf O_X\times_X M\longrightarrow M,\qquad(f,m)\longmapsto f\cdot m,
\]
使后者是一个 $X$-\Term{morphism}{态射}（即 $\pi(f\cdot m)=\pi(m)$），并且对每个 $x\in X(S)$，在 $M(x)=\{m\in M(S)\mid\pi(m)=x\}$ 上给出 $\mathbf O(S)$-\Term{module}{模}结构。
这时，对每个 $Y\in\operatorname{Ob}(\widehat{(\mathrm{Sch})}_{/X})$（不要求\Term{representable}{可表示}），$\operatorname{Hom}_X(Y,M)=\Gamma(M_Y/Y)$ 以关于 $Y$ \Term{functorial}{函子性}的方式成为 $\mathbf O(Y)$-\Term{module}{模}。

\SourcePage{29}
\textbf{4.4. \Term{diagonalizable-group}{可对角化群}。}

\textbf{4.4.1.} $\mathbb G_m$ 的构造可以推广如下。设 $M$ 是\Term{abelian-group}{交换群}，$M_{\mathbb Z}$ 是由 4.1 的方法赋予它的\Term{group-scheme}{群概形}。考虑由下式定义的\Term{functor}{函子}：
\[
\mathbf D(M)(S)=\operatorname{Hom}_{\mathrm{groupes}}(M,\mathbf G_m(S))\simeq\operatorname{Hom}_{S\text{-}\mathrm{Gr.}}(M_S,\mathbb G_{m,S}).
\]
这是一个交换的 $\widehat{(\mathrm{Sch})}$-\Term{group}{群}，可以由一个\Term{group-scheme}{群概形}表示，记为 $\mathrm D(M)$；因而按定义有
\[
\mathrm D(M)\simeq\underline{\operatorname{Hom}}_{(\mathrm{Sch})\text{-}\mathrm{Gr.}}(M_{\mathbb Z},\mathbb G_m).
\]
事实上，令
\[
\mathrm D(M)=\operatorname{Spec}\mathbb Z[M],
\]
其中 $\mathbb Z[M]$ 为\Term{group}{群} $M$ 在\Term{ring}{环} $\mathbb Z$ 上的\Term{group-algebra}{群代数}；由\Term{algebra}{代数} $\mathbb Z[M]$ 的定义即有
\[
\mathbf D(M)(S)=\operatorname{Hom}_{\mathrm{Alg.}}(\mathbb Z[M],\Gamma(S,\mathcal O_S))\simeq\operatorname{Hom}_{\mathrm{groupes}}(M,\Gamma(S,\mathcal O_S)^\times).
\]
\textbf{4.4.2.} 因此，对每个\Term{scheme}{概形} $S$，都有一个在 $S$ 上\Term{affine}{仿射}的 $S$-\Term{group-scheme}{群概形}
\[
\mathrm D_S(M)=\mathrm D(M)_S=\underline{\operatorname{Hom}}_{(\mathrm{Sch})\text{-}\mathrm{Gr.}}(M_{\mathbb Z},\mathbb G_m)_S=\underline{\operatorname{Hom}}_{S\text{-}\mathrm{Gr.}}(M_S,\mathbb G_{m,S}).
\]
它对应于 $\mathcal O_S$-\Term{bialgebra}{双代数} $\mathcal O_S[M]$，\Term{diagonal-map}{对角映射}与\Term{augmentation}{增广}定义为
\[
\Delta(x)=x\otimes x,\qquad\varepsilon(x)=1\quad(x\in M).
\]
\textbf{4.4.3.} 若 $f:M\to N$ 是\Term{abelian-group}{交换群}\Term{homomorphism}{同态}，就以显然的方式定义一个 $S$-\Term{group}{群}\Term{morphism}{态射}
\[
\mathrm D_S(f):\mathrm D_S(N)\longrightarrow\mathrm D_S(M),
\]
由此得到一个\Term{functor}{函子}
\[
\mathrm D_S:M\longmapsto\mathrm D_S(M),
\]
从\Term{abelian-group}{交换群}\Term{category}{范畴}到在 $S$ 上\Term{affine}{仿射}的 $S$-\Term{group}{群}\Term{category}{范畴}。也可以把它定义为\Term{functor}{函子} $M\mapsto M_S$ 与\Term{functor}{函子} $M_S\mapsto\underline{\operatorname{Hom}}_{S\text{-}\mathrm{Gr.}}(M_S,\mathbb G_{m,S})$ 的复合。此\Term{functor}{函子}与\Term{base-extension}{基扩张}交换。
与某个形如 $\mathrm D_S(M)$ 的\Term{group}{群}\Term{isomorphism}{同构}的 $S$-\Term{group}{群}称为\Term{diagonalizable}{可对角化}的。注意，$M$ 的元素可以解释为 $\mathrm D_S(M)$ 的某些\Term{character}{特征标}，即 $\operatorname{Hom}_{S\text{-}\mathrm{Gr.}}(\mathrm D_S(M),\mathbb G_{m,S})$ 的某些元素。（参见 VIII，1。）

\textbf{4.4.4.} 给出几个\Term{diagonalizable-group}{可对角化群}的例子。首先有
\[
\mathrm D(\mathbb Z)=\mathbb G_m,\qquad\mathrm D(\mathbb Z^r)=(\mathbb G_m)^r.
\]
记
\[
\mu_n=\mathrm D(\mathbb Z/n\mathbb Z),
\]
称为 $n$ 次\Term{roots-of-unity}{单位根群}。事实上，
\[
\mu_n(S)=\operatorname{Hom}_{\mathrm{groupes}}(\mathbb Z/n\mathbb Z,\Gamma(S,\mathcal O_S)^*)=\{f\in\Gamma(S,\mathcal O_S)\mid f^n=1\}.
\]
\SourcePage{30}
$S$-\Term{group}{群} $\mu_{n,S}$ 对应于 $\mathcal O_S$-\Term{algebra}{代数} $\mathcal O_S[T]/(T^n-1)$。特别地，设 $S$ 是一个\Term{characteristic}{特征} $p=n$ 的域 $k$ 的\Term{spectrum}{谱}。令 $T-1=s$，得到
\[
k[T]/(T^p-1)=k[s]/(s^p),
\]
这表明 $\mu_{p,k}$ 的底\Term{topological-space}{拓扑空间}只有一个点，该点的\Term{local-ring}{局部环}就是 \Term{artinian}{Artin} $k$-\Term{algebra}{代数} $k[s]/(s^p)$。（在同样的思路下，指出：$S$-\Term{scheme}{概形} $\mathbb G_{a,S}$、$\mathbb G_{m,S}$、$\mathbb O_S$ 在 $S$ 上\Term{smooth}{光滑}；$\mathrm D_S(M)$ 在 $S$ 上\Term{flat}{平坦}；它在 $S$ 上\Term{formally-smooth}{形式光滑}（分别为\Term{smooth}{光滑}），当且仅当 $S$ 的任何\Term{residue-characteristic}{剩余特征}均不整除 $M$ 的\Term{torsion}{挠性}（分别还要求 $M$ 为\Term{finite-type}{有限型}），参见第 VIII 讲，2.1。）

\textbf{4.5. 群的其他例子。} 上述方法适用于“\Term{classical-group}{典型群}”（\Term{general-linear-group}{一般线性群} $\mathrm{GL}_n$、\Term{symplectic-group}{辛群} $\mathrm{Sp}_n$、\Term{orthogonal-group}{正交群} $\mathrm O_n$ 等）。例如，把 $\mathrm{GL}_n$ 定义为下列\Term{functor}{函子} $\mathbf{GL}_n$ 的\Term{representing-object}{表示对象}：
\[
\mathbf{GL}_n(S)=\mathrm{GL}(n,\Gamma(S,\mathcal O_S))=\operatorname{Aut}_{\mathcal O_S}(\mathcal O_S^n).
\]
例如，可以将它构造为 $\operatorname{Spec}\mathbb Z[T_{ij}]$（$1\leq i,j\leq n$）中由函数 $f=\det((T_{ij})_{i,j=1}^n)$ 所定义的\Term{open-subset}{开集}，即 $\operatorname{Spec}\mathbb Z[T_{ij},f^{-1}]$。

\textbf{4.6. \Term{scheme}{概形}\Term{category}{范畴}中的模\Term{functor}{函子}。} 我们要对\Term{scheme}{概形} $S$ 上的每个 $\mathcal O_S$-\Term{module}{模}赋予一个 $\mathbf O_S$-\Term{module}{模}（$\mathbf O_S$ 表示 4.3.3 中引入的\Term{ring}{环}\Term{functor}{函子}）。这可以用两种不同的方式完成。确切地说：

\begin{SGAresult}{定义}{4.6.1}
设 $S$ 为\Term{scheme}{概形}。对每个 $\mathcal O_S$-\Term{module}{模} $\mathcal F$，用 $\mathbf V(\mathcal F)$ 与 $\mathbf W(\mathcal F)$ 表示由下式定义的 $(\mathrm{Sch})_{/S}$ 上的\Term{contravariant-functor}{反变函子}：
\[
\begin{aligned}
\mathbf V(\mathcal F)(S')&=\operatorname{Hom}_{\mathcal O_{S'}}(\mathcal F\otimes_{\mathcal O_S}\mathcal O_{S'},\mathcal O_{S'}),\\
\mathbf W(\mathcal F)(S')&=\Gamma(S',\mathcal F\otimes_{\mathcal O_S}\mathcal O_{S'}),
\end{aligned}
\]
其中 $\mathcal F\otimes_{\mathcal O_S}\mathcal O_{S'}$ 表示 $\mathcal O_S$-\Term{module}{模} $\mathcal F$ 在 $S'$ 上的\Term{inverse-image}{逆像}。
\end{SGAresult}
于是 $\mathbf V(\mathcal F)$ 与 $\mathbf W(\mathcal F)$ 以显然的方式带有 $\mathbf O_S$-\Term{module}{模}结构（回忆 $\mathbf O_S(S')=\Gamma(S',\mathcal O_{S'})=\mathbf W(\mathcal O_S)(S')$）；这样实际上定义了从 $\mathcal O_S$-\Term{module}{模}\Term{category}{范畴}到 $\mathbf O_S$-\Term{module}{模}\Term{category}{范畴}的两个\Term{functor}{函子} $\mathbf V$ 与 $\mathbf W$，前者\Term{contravariant}{反变}，后者\Term{covariant}{协变}。
本段以下仅考虑所涉及的 $\mathcal O_S$-\Term{module}{模}均为\Term{quasi-coherent}{拟凝聚}的情形；也就是说，把 $\mathbf V$ 和 $\mathbf W$ 看成从\Term{quasi-coherent}{拟凝聚} $\mathcal O_S$-\Term{module}{模}\Term{category}{范畴} $(\mathcal O_S\text{-}\mathrm{Mod.q.c.})$ 到 $\mathbf O_S$-\Term{module}{模}\Term{category}{范畴}的\Term{functor}{函子}：
\[
\begin{aligned}
\mathbf V &: (\mathcal O_S\text{-}\mathrm{Mod.q.c.})^\circ\longrightarrow(\mathbf O_S\text{-}\mathrm{Mod.}),\\
\mathbf W &: (\mathcal O_S\text{-}\mathrm{Mod.q.c.})\longrightarrow(\mathbf O_S\text{-}\mathrm{Mod.}).
\end{aligned}
\]
\SourcePage{31}
\begin{SGAresult}{注记}{4.6.1.1}
\footnote{编者注（31）：为便于后文引用，增添了编号 4.6.1.1。} 读者会注意到，下文中只涉及\Term{functor}{函子} $\mathbf W$ 的所有命题，对任意 $\mathcal O_S$-\Term{module}{模}都成立，无须要求\Term{quasi-coherent}{拟凝聚}。
\end{SGAresult}
\begin{SGAresult}{命题}{4.6.2}
（i）\Term{functor}{函子} $\mathbf V$ 与 $\mathbf W$ 同\Term{base-extension}{基扩张}交换：若 $S'$ 在 $S$ 上，且 $\mathcal F$ 是\Term{quasi-coherent}{拟凝聚} $\mathcal O_S$-\Term{module}{模}，则
\[
\mathbf V(\mathcal F\otimes\mathcal O_{S'})\simeq\mathbf V(\mathcal F)_{S'},\qquad\mathbf W(\mathcal F\otimes\mathcal O_{S'})\simeq\mathbf W(\mathcal F)_{S'}.
\]
（ii）\Term{functor}{函子} $\mathbf V$ 与 $\mathbf W$ 是\Term{fully-faithful}{全忠实}的：\Term{canonical}{典范}映射
\[
\begin{aligned}
\operatorname{Hom}_{\mathcal O_S}(\mathcal F,\mathcal F')&\longrightarrow\operatorname{Hom}_{\mathbf O_S}(\mathbf V(\mathcal F'),\mathbf V(\mathcal F)),\\
\operatorname{Hom}_{\mathcal O_S}(\mathcal F,\mathcal F')&\longrightarrow\operatorname{Hom}_{\mathbf O_S}(\mathbf W(\mathcal F),\mathbf W(\mathcal F'))
\end{aligned}
\]
是\Term{bijection}{双射}。
（iii）\Term{functor}{函子} $\mathbf V$ 与 $\mathbf W$ 是\Term{additive}{加性}的：
\[
\mathbf V(\mathcal F\oplus\mathcal F')\simeq\mathbf V(\mathcal F)\times_S\mathbf V(\mathcal F'),\qquad
\mathbf W(\mathcal F\oplus\mathcal F')\simeq\mathbf W(\mathcal F)\times_S\mathbf W(\mathcal F').
\]
\end{SGAresult}
\begin{proof}
（i）与（iii）由定义即明。对（ii），取 $S'$ 为 $S$ 的\Term{open-subset}{开集}。证明留给读者（对于 $\mathbf V$，使用 EGA II，1.7.14）。
\end{proof}
回忆（参见 3.1.4），若 $\mathrm F$、$\mathrm F'$ 是 $\mathbf O_S$-\Term{module}{模}，则 $\underline{\operatorname{Hom}}_{\mathbf O_S}(\mathrm F,\mathrm F')$ 表示向每个 $S'\to S$ 赋予 $\operatorname{Hom}_{\mathbf O_{S'}}(\mathrm F_{S'},\mathrm F'_{S'})$ 的 $S$-\Term{functor}{函子}（取值于\Term{abelian-group}{交换群}）。

\begin{SGAresult}{命题}{4.6.3}
\footnote{编者注（32）：增添了\Term{isomorphism}{同构} $\underline{\operatorname{Hom}}_{\mathbf O_S}(\mathbf W(\mathcal F),\mathbf W(\mathcal F'))\simeq\underline{\operatorname{Hom}}_{\mathbf O_S}(\mathbf V(\mathcal F'),\mathbf V(\mathcal F))$。} 在 $(\mathbf O_S\text{-}\mathrm{Mod.})$ 中有\Term{canonical}{典范}\Term{morphism}{态射}：
\[
\begin{tikzcd}[column sep=small]
\underline{\operatorname{Hom}}_{\mathbf O_S}(\mathbf W(\mathcal F),\mathbf W(\mathcal F'))\arrow[rr,equal]&&\underline{\operatorname{Hom}}_{\mathbf O_S}(\mathbf V(\mathcal F'),\mathbf V(\mathcal F))\\
&\mathbf W(\mathcal{H}om_{\mathcal O_S}(\mathcal F,\mathcal F'))\arrow[ul]\arrow[ur]&
\end{tikzcd}
\]
\end{SGAresult}
这立即由 4.6.2（i）与（ii）得到。

\textbf{记号 4.6.3.1。} 设 $\mathcal F$ 为\Term{quasi-coherent}{拟凝聚} $\mathcal O_S$-\Term{module}{模}。已知（EGA II，1.7.8），$S$-\Term{functor}{函子} $\mathbf V(\mathcal F)$ 可由一个在 $S$ 上\Term{affine}{仿射}的 $S$-\Term{scheme}{概形}表示；把它记为 $\mathbb V(\mathcal F)$，称为 $\mathcal F$ 所定义的\Term{vector-fibration}{向量纤维化}\footnote{编者注（33）：这里用“\Term{vector-fibration}{向量纤维化}”替换了“\Term{vector-bundle}{向量丛}”；按当前用法，“秩 $r$ 的\Term{vector-bundle}{向量丛}”指\Term{locally-trivial}{局部平凡}且秩为 $r$ 的\Term{vector-fibration}{向量纤维化}，即其\Term{section-sheaf}{截面层}\Term{local-scheme}{局部}\Term{isomorphism}{同构}于 $\mathcal O_S^{\oplus r}$。}：
\[
\mathbb V(\mathcal F)=\operatorname{Spec}(\mathrm S(\mathcal F)),
\]
\SourcePage{32}
其中 $\mathrm S(\mathcal F)$ 表示 $\mathcal O_S$-\Term{module}{模} $\mathcal F$ 的\Term{symmetric-algebra}{对称代数}。\footnote{编者注（34）：在此指出文献 [Ni04]（分别为 [Ni02]）：它们证明，若 $S$ 是 Noether \Term{scheme}{概形}，$\mathcal F$ 是\Term{coherent}{凝聚} $\mathcal O_S$-\Term{module}{模}，则 $\mathbf W(\mathcal F)$（分别为向每个 $T\to S$ 赋予 $\operatorname{Aut}_{\mathcal O_T}(\mathcal F\otimes\mathcal O_T)$ 的 $S$-\Term{group}{群}）\Term{representable}{可表示}，当且仅当 $\mathcal F$ \Term{locally-free}{局部自由}。}

\begin{SGAresult}{命题}{4.6.4}
设 $\mathcal F$、$\mathcal F'$ 为两个\Term{quasi-coherent}{拟凝聚} $\mathcal O_S$-\Term{module}{模}，$\mathcal A$ 为\Term{quasi-coherent}{拟凝聚} $\mathcal O_S$-\Term{algebra}{代数}。有\Term{functorial}{函子性}\Term{isomorphism}{同构}：
\[
\operatorname{Hom}_S\bigl(\operatorname{Spec}(\mathcal A),\underline{\operatorname{Hom}}_{\mathbf O_S}(\mathbf W(\mathcal F'),\mathbf W(\mathcal F))\bigr)\xrightarrow\sim\operatorname{Hom}_{\mathcal O_S}(\mathcal F',\mathcal F\otimes_{\mathcal O_S}\mathcal A).
\]
\end{SGAresult}
\begin{proof}
记 $X=\operatorname{Spec}(\mathcal A)$，则左端\Term{canonical}{典范}\Term{isomorphism}{同构}于 $\underline{\operatorname{Hom}}_{\mathbf O_S}(\mathbf W(\mathcal F'),\mathbf W(\mathcal F))(X)$；按定义，后者即
\[
\operatorname{Hom}_{\mathbf O_X}(\mathbf W(\mathcal F')_X,\mathbf W(\mathcal F)_X)\simeq\operatorname{Hom}_{\mathbf O_X}(\mathbf W(\mathcal F'\otimes\mathcal O_X),\mathbf W(\mathcal F\otimes\mathcal O_X))
\]
（参见 4.6.2（i））。由 4.6.2（ii），这也可写成
\[
\operatorname{Hom}_{\mathcal O_X}(\mathcal F'\otimes\mathcal O_X,\mathcal F\otimes\mathcal O_X)=\operatorname{Hom}_{\mathcal O_S}(\mathcal F',\pi_*(\pi^*(\mathcal F))),
\]
其中 $\pi:X\to S$ 表示\Term{structure-morphism}{结构态射}。由 EGA II，1.4.7，$\pi_*(\pi^*(\mathcal F))\simeq\mathcal F\otimes\mathcal A$，证明完成。
\end{proof}
\begin{SGAresult}{推论}{4.6.4.1}
有\Term{canonical}{典范}\Term{isomorphism}{同构}
\[
\mathbf W(\mathcal F\otimes\mathcal A)\simeq\underline{\operatorname{Hom}}_S(\operatorname{Spec}(\mathcal A),\mathbf W(\mathcal F)).
\]
\end{SGAresult}
\begin{proof}
\footnote{编者注（35）：以下对原文作了展开。} 设 $f:S'\to S$ 为 $S$-\Term{scheme}{概形}，且 $X'=X\times_S S'$，则有\Term{cartesian-square}{笛卡尔方块}
\[
\begin{tikzcd}
X'\arrow[r,"{f'}"]\arrow[d,"{\pi'}"']&X\arrow[d,"{\pi}"]\\
S'\arrow[r,"{f}"']&S.
\end{tikzcd}
\]
由 EGA II，1.5.2，$X'$ 在 $S'$ 上\Term{affine}{仿射}，且 $\pi'_*(\mathcal O_{X'})=f^*(\mathcal A)$。因而
\[
\underline{\operatorname{Hom}}_S(\operatorname{Spec}(\mathcal A),\mathbf W(\mathcal F))(S')=\operatorname{Hom}_{S'}\bigl(\operatorname{Spec}(f^*(\mathcal A)),\mathbf W(f^*(\mathcal F))\bigr).
\]
把 4.6.4 应用于 $f^*(\mathcal F)$、$\mathcal F'=\mathcal O_{S'}$ 与 $f^*(\mathcal A)$，可知这等于
\[
\Gamma(S',f^*(\mathcal F)\otimes f^*(\mathcal A))=\Gamma(S',f^*(\mathcal F\otimes\mathcal A))=\mathbf W(\mathcal F\otimes\mathcal A)(S').
\]
\end{proof}
\begin{SGAresult}{命题}{4.6.5}
若 $\mathcal F$、$\mathcal F'$ 为两个\Term{locally-free}{局部自由}\Term{finite-type}{有限型} $\mathcal O_S$-\Term{module}{模}，则 4.6.3 中的\Term{morphism}{态射}都是\Term{isomorphism}{同构}。
\end{SGAresult}
\begin{proof}
对每个 $S'\to S$，有
\[
\begin{aligned}
\mathbf W(\mathcal{H}om_{\mathcal O_S}(\mathcal F,\mathcal F'))(S')
&=\Gamma(S',\mathcal{H}om_{\mathcal O_S}(\mathcal F,\mathcal F')\otimes\mathcal O_{S'})\\
&=\operatorname{Hom}_{\mathcal O_{S'}}(\mathcal F\otimes\mathcal O_{S'},\mathcal F'\otimes\mathcal O_{S'}).
\end{aligned}
\]
由 4.6.2（i）与（ii），右端确实\Term{isomorphism}{同构}于 $\underline{\operatorname{Hom}}_{\mathbf O_S}(\mathbf W(\mathcal F),\mathbf W(\mathcal F'))(S')$ 及 $\underline{\operatorname{Hom}}_{\mathbf O_S}(\mathbf V(\mathcal F'),\mathbf V(\mathcal F))(S')$。
\end{proof}
\begin{SGAresult}{推论}{4.6.5.1}
设 $\mathcal F$ 为\Term{locally-free}{局部自由}\Term{finite-type}{有限型} $\mathcal O_S$-\Term{module}{模}。令 $\mathcal F^\vee=\mathcal{H}om_{\mathcal O_S}(\mathcal F,\mathcal O_S)$，则有\Term{canonical}{典范}\Term{isomorphism}{同构}：
\SourcePage{33}
\[
\begin{aligned}
\mathbf W(\mathcal F^\vee)&\simeq\underline{\operatorname{Hom}}_{\mathbf O_S}(\mathbf W(\mathcal F),\mathbf O_S)\simeq\mathbf V(\mathcal F),\\
\mathbf V(\mathcal F^\vee)&\simeq\underline{\operatorname{Hom}}_{\mathbf O_S}(\mathbf V(\mathcal F),\mathbf O_S)\simeq\mathbf W(\mathcal F).
\end{aligned}
\]
\end{SGAresult}
最后还有以下命题：
\begin{SGAresult}{命题}{4.6.6}
设 $f:\mathcal F\to\mathcal F'$ 为\Term{locally-free}{局部自由}\Term{finite-rank}{有限秩} $\mathcal O_S$-\Term{module-morphism}{模态射}。$\mathbf W(f):\mathbf W(\mathcal F)\to\mathbf W(\mathcal F')$ 为\Term{monomorphism}{单态射}，当且仅当 $f$ 在\Term{local-scheme}{局部}把 $\mathcal F$ 等同于 $\mathcal F'$ 的一个\Term{direct-summand}{直和因子}。
\end{SGAresult}
\begin{proof}
正向命题本质上已包含在 EGA $0_{\mathrm I}$，5.5.5 中。\footnote{编者注（36）：下一句作了展开。} 反之，若 $\mathcal F$ \Term{local-scheme}{局部}为 $\mathcal F'$ 的\Term{direct-summand}{直和因子}，则对每个 $\pi:S'\to S$，$\pi^*\mathcal F$ 是 $\pi^*\mathcal F'$ 的\Term{submodule}{子模}（因为\Term{local-scheme}{局部}为\Term{direct-summand}{直和因子}），故 $\mathbf W(\mathcal F)(S')=\Gamma(S',\pi^*\mathcal F)$ 是 $\mathbf W(\mathcal F')(S')=\Gamma(S',\pi^*\mathcal F')$ 的\Term{submodule}{子模}。
\end{proof}

\textbf{4.7. $G$-$\mathcal O_S$-\Term{module-category}{模范畴}。} 设 $G$ 是 $S$-\Term{group}{群}，$\mathcal F$ 是 $\mathcal O_S$-\Term{module}{模}；$\mathbf W(\mathcal F)$ 带有 $\mathbf O_S$-\Term{module}{模}结构。
\begin{SGAresult}{定义}{4.7.1}
$\mathcal F$ 上的 $G$-$\mathcal O_S$-\Term{module}{模}结构，是指 $\mathbf W(\mathcal F)$ 上的 $h_G$-$\mathbf O_S$-\Term{module}{模}结构（参见 3.2）。$G$-$\mathcal O_S$-\Term{module-morphism}{模态射}按定义为对应的 $h_G$-$\mathbf O_S$-\Term{module}{模}的\Term{morphism}{态射}。由此得到\Term{category}{范畴} $(G\text{-}\mathcal O_S\text{-}\mathrm{Mod.})$；用 $(G\text{-}\mathcal O_S\text{-}\mathrm{Mod.q.c.})$ 表示由作为 $\mathcal O_S$-\Term{module}{模}为\Term{quasi-coherent}{拟凝聚}的 $G$-$\mathcal O_S$-\Term{module}{模}组成的\Term{full-subcategory}{满子范畴}。
\end{SGAresult}
因而，由 3.2，在 $\mathcal F$ 上给定 $G$-$\mathcal O_S$-\Term{module}{模}结构，就是给定一个 $\widehat{(\mathrm{Sch})}_{/S}$-\Term{group}{群}\Term{morphism}{态射}
\[
\rho:h_G\longrightarrow\underline{\operatorname{Aut}}_{\mathbf O_S}(\mathbf W(\mathcal F)).
\]
\begin{SGAresult}{注记}{4.7.1.0}
\footnote{编者注（37）：此注记取自 [DG]，II，§2，1.1，为增补。} 由 4.6.2，有 $S$-\Term{group}{群}\Term{functor}{函子}的\Term{anti-isomorphism}{反同构}
\[
\underline{\operatorname{Aut}}_{\mathbf O_S}(\mathbf W(\mathcal F))\simeq\underline{\operatorname{Aut}}_{\mathbf O_S}(\mathbf V(\mathcal F)).\tag{$\dagger$}
\]
可见，在 $\mathbf W(\mathcal F)$ 上或在 $\mathbf V(\mathcal F)$ 上给定 $h_G$-$\mathbf O_S$-\Term{module}{模}结构是\Term{equivalence}{等价}的。事实上，设 $\rho:h_G\to\underline{\operatorname{Aut}}_{\mathbf O_S}(\mathbf W(\mathcal F))$ 如上。对每个 $T\to S$ 与 $g\in G(T)$，用 $\rho^*(g)$ 表示 $\rho(g)$ 在\Term{anti-isomorphism}{反同构}（$\dagger$）下的像；于是 $\rho^*(gh)=\rho^*(h)\circ\rho^*(g)$，即 $\rho^*$ 在 $\mathbf V(\mathcal F)$ 上定义一个“右”$h_G$-$\mathbf O_S$-\Term{module}{模}结构。令 $\rho^\vee(g)=\rho^*(g^{-1})$，便确实得到 $\mathbf V(\mathcal F)$ 上的 $h_G$-$\mathbf O_S$-\Term{module}{模}结构，且给定它\Term{equivalence}{等价}于给定 $\rho$。
\end{SGAresult}
\begin{SGAresult}{注记}{4.7.1.1}
可以说，刚才定义的\Term{category}{范畴}由以下\Term{cartesian-square}{笛卡尔方块}构造：
\[
\begin{tikzcd}[column sep=small]
(G\text{-}\mathcal O_S\text{-}\mathrm{Mod.q.c.})\arrow[r,hook]\arrow[d]&(G\text{-}\mathcal O_S\text{-}\mathrm{Mod.})\arrow[r]\arrow[d]&(h_G\text{-}\mathbf O_S\text{-}\mathrm{Mod.})\arrow[d,"{\text{遗忘}}"]\\
(\mathcal O_S\text{-}\mathrm{Mod.q.c.})\arrow[r,hook]&(\mathcal O_S\text{-}\mathrm{Mod.})\arrow[r,"{\mathbf W}"']&(\mathbf O_S\text{-}\mathrm{Mod.}).
\end{tikzcd}
\]
\SourcePage{34}
\footnote{编者注（38）：修正了原文，删去了\Term{category}{范畴} $(G\text{-}\mathcal O_S\text{-}\mathrm{Mod.})$ 是\Term{abelian-category}{阿贝尔范畴}的断言，参见下文 4.7.2.1。} \Term{category}{范畴} $(\mathcal O_S\text{-}\mathrm{Mod.})$ 与 $(\mathbf O_S\text{-}\mathrm{Mod.})$ 都是\Term{abelian-category}{阿贝尔范畴}，但须注意，\Term{functor}{函子} $\mathbf W$ 一般既非\Term{left-exact}{左正合}，也非\Term{right-exact}{右正合}。
\end{SGAresult}
\begin{SGAresult}{注记}{4.7.1.2}
\footnote{编者注（39）：此注记为增补。} 设 $\mathcal F$ 为 $G$-$\mathcal O_S$-\Term{module}{模}。\Term{invariants-subsheaf}{不变量子层} $\mathcal F^G$ 定义如下：对 $S$ 的每个\Term{open-subset}{开集} $U$，
\[
\mathcal F^G(U)=\mathbf W(\mathcal F)^G(U)=\{x\in\mathcal F(U)\mid g\cdot x_{S'}=x_{S'}\ \text{对所有 }S'\xrightarrow f U,\ g\in G(S')\},
\]
其中 $x_{S'}$ 表示 $x$ 在 $\Gamma(S',f^*(\mathcal F))=\Gamma(U,f_*f^*(\mathcal F))$ 中的像。
须注意，自然\Term{morphism}{态射} $\mathbf W(\mathcal F^G)\to\mathbf W(\mathcal F)^G$ 一般并非\Term{isomorphism}{同构}。例如，若 $S=\operatorname{Spec}(\mathbb Z)$，$G$ 为\Term{constant-group}{常群} $\mathbb Z/2\mathbb Z=\{1,\tau\}$，以 $\tau\cdot1=-1$ \Term{action}{作用}于 $\mathcal F=\mathcal O_S$，则 $\mathcal F^G=0$；但若 $R$ 为 $\mathbb F_2$-\Term{algebra}{代数}，则 $\mathbf W(\mathcal F)^G(\operatorname{Spec}(R))=R$。
\end{SGAresult}

\textbf{4.7.2.} 从此至 4.7 末，假设 $G$ 在 $S$ 上\Term{affine}{仿射}。\footnote{编者注（40）：关于把 4.7.2 的结果推广到 $G$ 不必\Term{affine}{仿射}、但 $G$ 与 $\mathcal F$ 都假设在 $S$ 上\Term{flat}{平坦}的情形，参见 $\mathrm{VI_B}$，§§11.1--11.6。} 于是由 4.6.4，给定一个 $S$-\Term{morphism-of-functors}{函子态射}
\[
\rho:h_G\longrightarrow\underline{\operatorname{End}}_{\mathbf O_S}(\mathbf W(\mathcal F))
\]
\Term{equivalence}{等价}于给定一个 $\mathcal O_S$-\Term{module-morphism}{模态射}
\[
\mu:\mathcal F\longrightarrow\mathcal F\otimes_{\mathcal O_S}\mathcal A(G).
\]
说 $\rho$ 是 $\widehat{(\mathrm{Sch})}_{/S}$-\Term{group}{群}\Term{morphism}{态射}，便\Term{equivalence}{等价}于说 $\mu$ 满足以下公理：

（CM 1）下图交换：
\[
\begin{tikzcd}
\mathcal F\arrow[r,"{\mu}"]\arrow[d,"{\mu}"']&\mathcal F\otimes_{\mathcal O_S}\mathcal A(G)\arrow[d,"{\mathrm{id}\otimes\Delta}"]\\
\mathcal F\otimes_{\mathcal O_S}\mathcal A(G)\arrow[r,"{\mu\otimes\mathrm{id}}"']&\mathcal F\otimes_{\mathcal O_S}\mathcal A(G)\otimes_{\mathcal O_S}\mathcal A(G).
\end{tikzcd}
\]
（CM 2）以下复合为恒等：
\[
\mathcal F\xrightarrow\mu\mathcal F\otimes\mathcal A(G)\xrightarrow{\mathrm{id}\otimes\varepsilon}\mathcal F\otimes\mathcal O_S\xrightarrow\sim\mathcal F.
\]
公理（CM 1）与（CM 2）就是\Term{bialgebra}{双代数} $\mathcal A(G)$ 上的（右）\Term{comodule}{余模}结构的公理。\footnote{编者注（41）：左 $G$-$\mathcal O_S$-\Term{module}{模}自然对应于右 $\mathcal A(G)$-\Term{comodule}{余模}。}

\SourcePage{35}
令 $\mathcal A=\mathcal A(G)$。若 $\mathcal F$ 与 $\mathcal F'$ 为 $\mathcal A$-\Term{comodule}{余模}，则\Term{comodule}{余模}\Term{morphism}{态射} $f:\mathcal F\to\mathcal F'$ 是使下图交换的 $\mathcal O_S$-\Term{module-morphism}{模态射}：
\[
\begin{tikzcd}
\mathcal F\arrow[r,"{f}"]\arrow[d,"{\mu_{\mathcal F}}"']&\mathcal F'\arrow[d,"{\mu_{\mathcal F'}}"]\\
\mathcal F\otimes\mathcal A\arrow[r,"{f\otimes\mathrm{id}}"']&\mathcal F'\otimes\mathcal A.
\end{tikzcd}
\]
由此得到\Term{category}{范畴} $(\mathcal A\text{-}\mathrm{Comod.})$；用 $(\mathcal A\text{-}\mathrm{Comod.q.c.})$ 表示由作为 $\mathcal O_S$-\Term{module}{模}为\Term{quasi-coherent}{拟凝聚}的 $\mathcal A$-\Term{comodule}{余模}组成的\Term{full-subcategory}{满子范畴}。因而得到了：
\begin{SGAresult}{命题}{4.7.2}
设 $G$ 为在 $S$ 上\Term{affine}{仿射}的 $S$-\Term{group}{群}。有\Term{equivalence-of-categories}{范畴等价}：
\[
\begin{aligned}
(G\text{-}\mathcal O_S\text{-}\mathrm{Mod.})&\simeq(\mathcal A(G)\text{-}\mathrm{Comod.}),\\
(G\text{-}\mathcal O_S\text{-}\mathrm{Mod.q.c.})&\simeq(\mathcal A(G)\text{-}\mathrm{Comod.q.c.}).
\end{aligned}
\]
\footnote{编者注（42）：下一句为增补。} 若还有 $S=\operatorname{Spec}(\Lambda)$ \Term{affine}{仿射}，并记 $\Lambda[G]=\Gamma(S,\mathcal A(G))$，则有\Term{equivalence-of-categories}{范畴等价}
\[
(\mathcal A(G)\text{-}\mathrm{Comod.q.c.})\simeq(\Lambda[G]\text{-}\mathrm{Comod.}).
\]
\end{SGAresult}
\footnote{编者注（43）：下一段取自 [Se68，§1.3]，为增补。} 再假设 $\mathcal A=\mathcal A(G)$ 为\Term{flat}{平坦} $\mathcal O_S$-\Term{module}{模}。设 $\mathcal E$ 为 $\mathcal A$-\Term{comodule}{余模}，$\mathcal F$ 为 $\mathcal E$ 的子 $\mathcal O_S$-\Term{module}{模}。由于 $\mathcal A$ 在 $\mathcal O_S$ 上\Term{flat}{平坦}，可以把 $\mathcal F\otimes\mathcal A$（分别为 $\mathcal F\otimes\mathcal A\otimes\mathcal A$）等同于 $\mathcal E\otimes\mathcal A$（分别为 $\mathcal E\otimes\mathcal A\otimes\mathcal A$）的子 $\mathcal O_S$-\Term{module}{模}。假设 $\mu_{\mathcal E}$ 把 $\mathcal F$ 映入 $\mathcal F\otimes\mathcal A$，则其\Term{restriction}{限制} $\mu_{\mathcal F}:\mathcal F\to\mathcal F\otimes\mathcal A$ 在 $\mathcal F$ 上给出 $\mathcal A$-\Term{comodule}{余模}结构；称 $\mathcal F$ 为 $\mathcal E$ 的\Term{subcomodule}{子余模}。取商后，$\mu_{\mathcal E}$ 定义一个 $\mathcal O_S$-\Term{module-morphism}{模态射} $\mathcal E/\mathcal F\to(\mathcal E/\mathcal F)\otimes\mathcal A$，从而在 $\mathcal E/\mathcal F$ 上给出 $\mathcal A$-\Term{comodule}{余模}结构。若 $f:\mathcal E\to\mathcal E'$ 为 $\mathcal A$-\Term{comodule}{余模}\Term{morphism}{态射}，则 $\operatorname{Ker}f$（分别为 $\operatorname{Im}f$）是 $\mathcal E$（分别为 $\mathcal E'$）的子 $\mathcal A$-\Term{comodule}{余模}，且 $f$ 诱导一个 $\mathcal A$-\Term{comodule}{余模}\Term{isomorphism}{同构}：$\mathcal E/\operatorname{Ker}f\xrightarrow\sim\operatorname{Im}f$。此外，若 $\mathcal E$ 与 $\mathcal E'$ 为\Term{quasi-coherent}{拟凝聚} $\mathcal O_S$-\Term{module}{模}，则 $\operatorname{Ker}f$ 与 $\operatorname{Im}f$ 也如此。因此，$(\mathcal A\text{-}\mathrm{Comod.})$ 与 $(\mathcal A\text{-}\mathrm{Comod.q.c.})$ 都是\Term{abelian-category}{阿贝尔范畴}。
\begin{SGAresult}{推论}{4.7.2.1}
假设 $G$ 在 $S$ 上\Term{affine}{仿射}且\Term{flat}{平坦}。则\Term{category}{范畴} $(G\text{-}\mathcal O_S\text{-}\mathrm{Mod.q.c.})$（分别为 $(G\text{-}\mathcal O_S\text{-}\mathrm{Mod.})$）是\Term{abelian-category}{阿贝尔范畴}，它\Term{equivalence}{等价}于在 $\mathcal O_S$ 上\Term{quasi-coherent}{拟凝聚}的 $\mathcal A(G)$-\Term{comodule}{余模}\Term{category}{范畴}（分别为 $\mathcal A(G)$-\Term{comodule}{余模}\Term{category}{范畴}）。
\end{SGAresult}

\textbf{4.7.3.} 现在假设 $G$ 为\Term{diagonalizable-group}{可对角化群}，也就是说，$\mathcal A(G)$ 是一个\Term{abelian-group}{交换群} $M$ 在\Term{ringed-sheaf}{环层} $\mathcal O_S$ 上的\Term{group-algebra}{群代数}。若 $\mathcal F$ 为 $\mathcal O_S$-\Term{module}{模}，则
\[
\mathcal F\otimes\mathcal A(G)=\coprod_{m\in M}\mathcal F\otimes m\mathcal O_S.
\]
\SourcePage{36}
因此，给定一个 $\mathcal O_S$-\Term{module-morphism}{模态射}
\[
\mu:\mathcal F\longrightarrow\mathcal F\otimes\mathcal A(G)
\]
\Term{equivalence}{等价}于给定 $\mathcal F$ 的一族 $\mathcal O_S$-\Term{endomorphism}{自同态} $(\mu_m)_{m\in M}$，使得对 $S$ 的一个\Term{open-subset}{开集}上 $\mathcal F$ 的每个\Term{section}{截面} $x$，$(\mu_m(x))$ 是\Term{direct-sum}{直和} $\coprod_{m\in M}\mathcal F$ 的一个\Term{section}{截面}（也就是说，在每个足够小的\Term{open-subset}{开集}上，只有\Term{finite}{有限}多个 $\mu_m(x)$ 的\Term{restriction}{限制}非零）。
由
\[
\mu(x)=\sum_{m\in M}\mu_m(x)\otimes m
\]
所定义的 $\mu$ 满足（CM 1）（分别为（CM 2）），当且仅当
\[
\mu_m\circ\mu_{m'}=\delta_{mm'}\mu_m,\qquad\left(\text{分别为 }\sum_{m\in M}\mu_m=\mathrm{Id}_{\mathcal F}\right),
\]
即 $\mu_m$ 是两两\Term{orthogonal}{正交}且和为恒等的\Term{projection-operator}{投影算子}。因此证明了：
\begin{SGAresult}{命题}{4.7.3}
若 $G=\mathrm D_S(M)$ 是\Term{diagonalizable}{可对角化} $S$-\Term{group}{群}，则\Term{quasi-coherent}{拟凝聚} $G$-$\mathcal O_S$-\Term{module}{模}\Term{category}{范畴}（分别为 $G$-$\mathcal O_S$-\Term{module}{模}\Term{category}{范畴}）\Term{equivalence}{等价}于 $M$ 型\Term{grading}{分次}的\Term{quasi-coherent}{拟凝聚} $\mathcal O_S$-\Term{module}{模}\Term{category}{范畴}（分别为 $\mathcal O_S$-\Term{module}{模}\Term{category}{范畴}）。
\end{SGAresult}
\textbf{注记。} 若 $\mathcal F$ 带有被 $G$ 的\Term{action}{作用}保持的 $\mathcal O_S$-\Term{algebraic-structure}{代数结构}，则 $\mathcal F$ 的\Term{grading}{分次}是\Term{algebra}{代数}的\Term{grading}{分次}。更确切地说：
\begin{SGAresult}{推论}{4.7.3.1}
\Term{functor}{函子} $\mathcal A\mapsto\operatorname{Spec}\mathcal A$ 在 $M$ 型\Term{grading}{分次}的\Term{quasi-coherent}{拟凝聚} $\mathcal O_S$-\Term{algebra}{代数}\Term{category}{范畴}与带有\Term{operator}{算子} $S$-\Term{group}{群} $G=\mathrm D_S(M)$、且在 $S$ 上\Term{affine}{仿射}的 $S$-\Term{scheme}{概形}\Term{category}{范畴}的\Term{opposite-category}{对偶范畴}之间诱导\Term{equivalence}{等价}。
\end{SGAresult}
\begin{SGAresult}{命题}{4.7.4}
设 $G$ 为\Term{diagonalizable}{可对角化} $S$-\Term{group}{群}。若 $0\to\mathcal F_1\to\mathcal F_2\to\mathcal F_3\to0$ 是\Term{quasi-coherent}{拟凝聚} $G$-$\mathcal O_S$-\Term{module}{模}的\Term{exact-sequence}{正合列}，并且作为 $\mathcal O_S$-\Term{module}{模}的序列\Term{split}{分裂}，则它作为 $G$-$\mathcal O_S$-\Term{module}{模}的序列也\Term{split}{分裂}。
\end{SGAresult}
\begin{proof}
若 $G=\mathrm D_S(M)$，则每个 $\mathcal F_i$ 由 $(\mathcal F_i)_m$ \Term{grading}{分次}，且对每个 $m\in M$，$\mathcal O_S$-\Term{module}{模}序列
\[
0\longrightarrow(\mathcal F_1)_m\longrightarrow(\mathcal F_2)_m\longrightarrow(\mathcal F_3)_m\longrightarrow0
\]
\Term{split}{分裂}。由前一命题即得结论。
\end{proof}

\SGAsection{5}{群的上同调}
\textbf{5.1. \Term{standard-complex}{标准复形}。}\footnote{编者注（44）：这个\Term{complex}{复形}常称为“Hochschild \Term{complex}{复形}”；例如参见 [DG70] 的 §II.3。} 设 $C$ 为\Term{category}{范畴}，$\mathbf G$ 为 $\widehat C$-\Term{group}{群}，$\mathbf O$ 为 $\widehat C$-\Term{ring}{环}，$\mathrm F$ 为 $\mathbf G$-$\mathbf O$-\Term{module}{模}。对 $n\geq0$，令
\[
C^n(\mathbf G,\mathrm F)=\operatorname{Hom}(\mathbf G^n,\mathrm F),\qquad
\underline C^n(\mathbf G,\mathrm F)=\underline{\operatorname{Hom}}(\mathbf G^n,\mathrm F),
\]
\SourcePage{37}
其中 $\mathbf G^0$ 是\Term{final-object}{终对象} $e$。于是 $\underline C^n(\mathbf G,\mathrm F)$（分别为 $C^n(\mathbf G,\mathrm F)$）以显然的方式带有 $\mathbf O$-\Term{module}{模}结构（分别为 $\Gamma(\mathbf O)$-\Term{module}{模}结构），且有
\[
C^n(\mathbf G,\mathrm F)\simeq\Gamma(\underline C^n(\mathbf G,\mathrm F)),\qquad
\underline C^n(\mathbf G,\mathrm F)(S)=C^n(\mathbf G_S,\mathrm F_S).
\]
给定 $C^n(\mathbf G,\mathrm F)$ 的一个元素，就是对每个 $S\in\operatorname{Ob}(C)$，以关于 $S$ \Term{functorial}{函子性}的方式给定 $\mathbf G(S)$ 到 $\mathrm F(S)$ 的一个 $n$-\Term{cochain}{上链}。回忆，\Term{boundary-operator}{边界算子}
\[
\partial:C^n(\mathbf G(S),\mathrm F(S))\longrightarrow C^{n+1}(\mathbf G(S),\mathrm F(S))
\]
由以下公式给出：
\[
\begin{aligned}
\partial f(g_1,\ldots,g_{n+1})={}&g_1f(g_2,\ldots,g_{n+1})\\
&+\sum_{i=1}^n(-1)^if(g_1,\ldots,g_ig_{i+1},\ldots,g_{n+1})\\
&+(-1)^{n+1}f(g_1,\ldots,g_n).
\end{aligned}
\]
它关于 $S$ 是\Term{functorial}{函子性}的，因而定义\Term{homomorphism}{同态}
\[
\partial:C^n(\mathbf G,\mathrm F)\longrightarrow C^{n+1}(\mathbf G,\mathrm F),
\]
满足 $\partial\circ\partial=0$。这样就定义了一个\Term{abelian-group}{交换群}的\Term{complex}{复形}（甚至是 $\Gamma(\mathbf O)$-\Term{module}{模}的\Term{complex}{复形}），记为 $C^*(\mathbf G,\mathrm F)$。同样定义 $\mathbf O$-\Term{module}{模}\Term{complex}{复形} $\underline C^*(\mathbf G,\mathrm F)$，并有
\[
C^*(\mathbf G,\mathrm F)=\Gamma(\underline C^*(\mathbf G,\mathrm F)).
\]
以 $H^n(\mathbf G,\mathrm F)$（分别为 $\underline H^n(\mathbf G,\mathrm F)$）表示\Term{complex}{复形} $C^*(\mathbf G,\mathrm F)$（分别为 $\underline C^*(\mathbf G,\mathrm F)$）的\Term{homology-group}{同调群}（分别为 $\widehat C$-\Term{homology-group}{同调群}）。特别地，
\[
\underline H^0(\mathbf G,\mathrm F)=\mathrm F^{\mathbf G},\qquad H^0(\mathbf G,\mathrm F)=\Gamma(\mathrm F^{\mathbf G}).
\]
\begin{SGAresult}{注记}{5.1.1}
\footnote{编者注（45）：此注记为增补。} 上述关于 $\partial$ 的“\Term{set}{集合}式”描述，便于核实 $\partial\circ\partial=0$。也可以用\Term{multiplication}{乘法} $m:\mathbf G\times\mathbf G\to\mathbf G$ 和\Term{action}{作用} $\mu:\mathbf G\times\mathrm F\to\mathrm F$ 来定义 $\partial$：对每个 $f\in C^n(\mathbf G,\mathrm F)$，
\[
\partial f=\mu\circ(\mathrm{id}_{\mathbf G}\times f)+\sum_{i=1}^n(-1)^if\circ(\mathrm{id}_{\mathbf G^{i-1}}\times m\times\mathrm{id}_{\mathbf G^{n-i}})+(-1)^{n+1}f\circ\operatorname{pr}_{[1,n]},
\]
其中 $\operatorname{pr}_{[1,n]}$ 表示从 $\mathbf G^{n+1}=\mathbf G^n\times\mathbf G$ 到 $\mathbf G^n$ 的\Term{projection}{投影}。同样，对每个 $S\in\operatorname{Ob}(C)$ 和 $f\in\underline C^n(\mathbf G,\mathrm F)(S)=C^n(\mathbf G_S,\mathrm F_S)$，有
\[
\partial f=\mu_S\circ(\mathrm{id}_{\mathbf G_S}\times f)+\sum_{i=1}^n(-1)^if\circ(\mathrm{id}_{\mathbf G_S^{i-1}}\times m_S\times\mathrm{id}_{\mathbf G_S^{n-i}})+(-1)^{n+1}f\circ\operatorname{pr}_{[1,n]},
\]
其中 $m_S$ 与 $\mu_S$ 由 $m$ 与 $\mu$ 作\Term{base-change}{基变换}得到。
\end{SGAresult}

\SourcePage{38}
\textbf{5.2.}\footnote{编者注（46）：以下回顾为增补。} 回忆（参见 §3），$(\mathbf G\text{-}\mathbf O\text{-}\mathrm{Mod.})$ 带有“逐个自变量”定义的\Term{abelian-category}{阿贝尔范畴}结构；因此，
\[
0\longrightarrow\mathrm F'\longrightarrow\mathrm F\longrightarrow\mathrm F''\longrightarrow0
\]
是 $\mathbf G$-$\mathbf O$-\Term{module}{模}的\Term{exact-sequence}{正合列}，当且仅当对每个 $S\in\operatorname{Ob}(C)$，\Term{abelian-group}{交换群}序列
\[
0\longrightarrow\mathrm F'(S)\longrightarrow\mathrm F(S)\longrightarrow\mathrm F''(S)\longrightarrow0
\]
\Term{exact}{正合}。
\footnote{编者注（47）：下一句为增补。} 假设 $C$ 为\Term{small-category}{小范畴}；由 3.2.1，$(\mathbf G\text{-}\mathbf O\text{-}\mathrm{Mod.})$ 有足够多的\Term{injective-object}{内射对象}，从而\Term{left-exact}{左正合}\Term{functor}{函子} $\underline H^0$ 与 $H^0$ 的\Term{derived-functor}{导出函子}有定义。下面要证明，\Term{functor}{函子} $\underline H^n$（分别为 $H^n$）正是 $\underline H^0$（分别为 $H^0$）的\Term{derived-functor}{导出函子}。

\begin{SGAresult}{定义}{5.2.0}
\footnote{编者注（48）：原文作了修改，以引入 5.2.0.1 与 5.2.0.2，它们将在\Term{theorem}{定理} 5.3.1 的证明中使用。} 对每个 $\mathbf O$-\Term{module}{模} $\mathrm P$，用 $\mathrm E(\mathrm P)$ 表示 $\widehat C$ 的\Term{object}{对象} $\underline{\operatorname{Hom}}(\mathbf G,\mathrm P)$，并赋予如下定义的 $\mathbf G$-$\mathbf O$-\Term{module}{模}结构：对每个 $S\in\operatorname{Ob}(C)$，有 $\underline{\operatorname{Hom}}(\mathbf G,\mathrm P)(S)=\operatorname{Hom}_S(\mathbf G_S,\mathrm P_S)$；令 $g\in\mathbf G(S)$ 与 $a\in\mathbf O(S)$ 按公式
\[
(g\cdot\varphi)(h)=\varphi(hg),\qquad(a\cdot\varphi)(h)=a\varphi(h)
\]
\Term{action}{作用}于 $\varphi\in\operatorname{Hom}_S(\mathbf G_S,\mathrm P_S)$，这里 $h\in\mathbf G(S')$、$S'\to S$ 任意。此外，对每个 $\varphi\in\operatorname{Hom}_S(\mathbf G_S,\mathrm P_S)$，令
\[
\varepsilon(\varphi)=\varphi(1)\in\mathrm P(S),
\]
其中 $1$ 表示 $\mathbf G(S)$ 的\Term{identity-element}{单位元}。
\end{SGAresult}
这定义了一个\Term{functor}{函子} $\mathrm E:(\mathbf O\text{-}\mathrm{Mod.})\to(\mathbf G\text{-}\mathbf O\text{-}\mathrm{Mod.})$ 与一个\Term{natural-transformation}{自然变换} $\varepsilon:\mathrm E\to\mathrm{Id}$，其中 $\mathrm{Id}$ 表示 $(\mathbf O\text{-}\mathrm{Mod.})$ 的\Term{identity-functor}{恒等函子}。

\begin{SGAresult}{注记}{5.2.0.1}
\footnote{编者注（48）：原文作了修改，以引入 5.2.0.1 与 5.2.0.2，它们将在\Term{theorem}{定理} 5.3.1 的证明中使用。} 下文中用 $\mathbf G_1$ 与 $\mathbf G_2$ 表示 $\mathbf G$ 的两个副本。则\Term{morphism}{态射}
\[
\mathbf G_1\times\mathrm E(\mathrm P)\longrightarrow\mathrm E(\mathrm P),\qquad(g_1,\varphi)\longmapsto\bigl(g_2\longmapsto\varphi(g_2g_1)\bigr)
\]
通过\Term{isomorphism}{同构}
\[
\begin{aligned}
\operatorname{Hom}(\mathbf G_1\times\mathrm E(\mathrm P),\mathrm E(\mathrm P))
&\simeq\operatorname{Hom}(\mathrm E(\mathrm P),\underline{\operatorname{Hom}}(\mathbf G_1,\underline{\operatorname{Hom}}(\mathbf G_2,\mathrm P)))\\
&\simeq\operatorname{Hom}(\mathrm E(\mathrm P),\underline{\operatorname{Hom}}(\mathbf G_2\times\mathbf G_1,\mathrm P))
\end{aligned}
\]
对应于\Term{morphism}{态射} $\varphi\mapsto((g_2,g_1)\mapsto\varphi(g_2g_1))$，即由\Term{multiplication}{乘法} $\mu_{\mathbf G}:\mathbf G\times\mathbf G\to\mathbf G$、$(g_2,g_1)\mapsto g_2g_1$ 诱导的\Term{morphism}{态射}
\[
\underline{\operatorname{Hom}}(\mathbf G,\mathrm P)\longrightarrow\underline{\operatorname{Hom}}(\mathbf G_2\times\mathbf G_1,\mathrm P).
\]
\end{SGAresult}

\SourcePage{39}
\begin{SGAresult}{\Term{lemma}{引理}}{5.2.0.2}
\footnote{编者注（48）：参见定义 5.2.0 的编者注；原书在此重复使用同一脚注号。} （i）\Term{functor}{函子} $\mathrm E$ 是\Term{forgetful-functor}{遗忘函子} $(\mathbf G\text{-}\mathbf O\text{-}\mathrm{Mod.})\to(\mathbf O\text{-}\mathrm{Mod.})$ 的\Term{right-adjoint}{右伴随}；更确切地说，对每个 $\mathrm M\in(\mathbf G\text{-}\mathbf O\text{-}\mathrm{Mod.})$ 与 $\mathrm P\in\operatorname{Ob}(\mathbf O\text{-}\mathrm{Mod.})$，$\varepsilon:\mathrm E\to\mathrm{Id}$ 诱导一个关于 $\mathrm M$、$\mathrm P$ \Term{functorial}{函子性}的\Term{bijection}{双射}
\[
\operatorname{Hom}_{\mathbf G\text{-}\mathbf O\text{-}\mathrm{Mod.}}(\mathrm M,\mathrm E(\mathrm P))\xrightarrow\sim\operatorname{Hom}_{\mathbf O\text{-}\mathrm{Mod.}}(\mathrm M,\mathrm P).
\]
（ii）因此，若 $\mathrm I$ 为 $(\mathbf O\text{-}\mathrm{Mod.})$ 的\Term{injective-object}{内射对象}，则 $\mathrm E(\mathrm I)$ 为 $(\mathbf G\text{-}\mathbf O\text{-}\mathrm{Mod.})$ 的\Term{injective-object}{内射对象}。
\end{SGAresult}
\begin{proof}
对每个 $\mathbf O$-\Term{morphism}{态射} $f:\mathrm M\to\mathrm P$，对应以下定义的 $\operatorname{Hom}_{\mathbf O}(\mathrm M,\mathrm E(\mathrm P))$ 的元素 $\varphi_f$。对每个 $S\in\operatorname{Ob}(C)$ 与 $m\in\mathrm M(S)$，$\varphi_f(m)$ 是如下定义的 $\operatorname{Hom}_S(\mathbf G_S,\mathrm P_S)$ 的元素：对每个 $g\in\mathbf G(S')$、$S'\to S$，
\[
\varphi_f(m)(g)=f(gm)\in\mathrm P(S').
\]
于是，对每个 $h\in\mathbf G(S)$，有 $\varphi_f(hm)=h\cdot f(m)$，即 $\varphi_f$ 为 $\operatorname{Hom}_{\mathbf G\text{-}\mathbf O\text{-}\mathrm{Mod.}}(\mathrm M,\mathrm E(\mathrm P))$ 的元素。
若 $\varphi\in\operatorname{Hom}_{\mathbf G\text{-}\mathbf O\text{-}\mathrm{Mod.}}(\mathrm M,\mathrm E(\mathrm P))$，并对每个 $m\in\mathrm M(S)$ 记 $f(m)=\varphi(m)(1)$，则
\[
\varphi_f(m)(g)=f(gm)=\varphi(gm)(1)=(g\cdot\varphi(m))(1)=\varphi(m)(g),
\]
即 $\varphi_f=\varphi$。反之，显然 $\varphi_f(m)(1)=f(m)$。这证明了（i），（ii）随即得出。
\end{proof}
\begin{SGAresult}{定义}{5.2.0.3}
设 $\mathrm M$ 为 $\mathbf G$-$\mathbf O$-\Term{module}{模}；$\mathrm M$ 的\Term{identity-map}{恒等映射}（把它看成 $\mathbf O$-\Term{module}{模}）通过\Term{adjunction}{伴随}对应于 $\mathbf G$-$\mathbf O$-\Term{module-morphism}{模态射}
\[
\mu_{\mathrm M}:\mathrm M\longrightarrow\mathrm E(\mathrm M),
\]
它满足：对每个 $S\in\operatorname{Ob}(C)$ 与 $m\in\mathrm M(S)$，$\mu_{\mathrm M}(m)$ 是由下式定义的\Term{morphism}{态射} $\mathbf G_S\to\mathrm M_S$：对每个 $S'\to S$ 和 $g\in\mathbf G(S')$，$\mu_{\mathrm M}(m)(g)=g\cdot m_{S'}\in\mathrm M(S')$。
\end{SGAresult}
注意 $\mu_{\mathrm M}$ 是\Term{monomorphism}{单态射}。事实上，$\varepsilon_{\mathrm M}:\mathrm E(\mathrm M)\to\mathrm M$ 是 $\mathbf O$-\Term{module-morphism}{模态射}，满足 $\varepsilon_{\mathrm M}\circ\mu_{\mathrm M}=\mathrm{id}_{\mathrm M}$；因此，作为 $\mathbf O$-\Term{module}{模}，$\mathrm M$ 是 $\mathrm E(\mathrm M)$ 的\Term{direct-summand}{直和因子}。

\begin{SGAresult}{命题}{5.2.1}
假设 $C$ 是\Term{small-category}{小范畴}，其中存在\Term{finite-product}{有限积}，且 $\mathbf G$ \Term{representable}{可表示}。则\Term{functor}{函子} $\underline H^n(\mathbf G,\ )$（分别为 $H^n(\mathbf G,\ )$）是在 $\mathbf G$-$\mathbf O$-\Term{module}{模}\Term{category}{范畴}上\Term{left-exact}{左正合}\Term{functor}{函子} $\underline H^0(\mathbf G,\ )$（分别为 $H^0(\mathbf G,\ )$）的\Term{derived-functor}{导出函子}。
\end{SGAresult}
\begin{proof}
根据熟知的一般结果\footnote{编者注（49）：参见 [Gr57]，2.2.1 与 2.3。此外，以下对原文作了展开。}，只须核实 $\underline H^n(\mathbf G,\ )$（分别为 $H^n(\mathbf G,\ )$）构成一个在正\Term{degree}{次数}\Term{effaceable}{可抹除}的\Term{cohomological-functor}{上同调函子}。
设
\[
0\longrightarrow\mathrm F'\longrightarrow\mathrm F\longrightarrow\mathrm F''\longrightarrow0
\]
是 $\mathbf G$-$\mathbf O$-\Term{module}{模}的\Term{exact-sequence}{正合列}，$S\in\operatorname{Ob}(C)$。依假设，$\mathbf G$ 可由一个\Term{object}{对象} $G\in\operatorname{Ob}(C)$ 表示，且 $C$ 中存在\Term{finite-product}{有限积}；特别地，
\SourcePage{40}
$C$ 有\Term{final-object}{终对象} $e$。因此，每个 $\mathbf G^n\times h_S$ 都可由 $G^n\times S$ 表示（其中 $G^0=e$），且序列
\[
0\longrightarrow\mathrm F'(G^n\times S)\longrightarrow\mathrm F(G^n\times S)\longrightarrow\mathrm F''(G^n\times S)\longrightarrow0
\]
\Term{exact}{正合}。这表明 $\mathbf O$-\Term{module}{模}序列
\[
0\longrightarrow\underline C^n(h_G,\mathrm F')\longrightarrow\underline C^n(h_G,\mathrm F)\longrightarrow\underline C^n(h_G,\mathrm F'')\longrightarrow0
\]
\Term{exact}{正合}。因此，把 $\underline C^*(\mathbf G,\ )$ 看作 $(\mathbf G\text{-}\mathbf O\text{-}\mathrm{Mod.})$ 上取值于 $(\mathbf O\text{-}\mathrm{Mod.})$ 的\Term{complex}{复形}\Term{category}{范畴}的\Term{functor}{函子}，它是\Term{exact}{正合}的。这表明 $\underline H^n(\mathbf G,\ )$ 确实构成\Term{cohomological-functor}{上同调函子}。由于\Term{functor}{函子} $\Gamma$ \Term{exact}{正合}，$H^n(\mathbf G,\ )$ 也如此。现在只须证明下列\Term{lemma}{引理}。
\end{proof}
\begin{SGAresult}{\Term{lemma}{引理}}{5.2.2}
对每个 $\mathrm P\in\operatorname{Ob}(\mathbf O\text{-}\mathrm{Mod.})$，有
\[
\underline H^n(\mathbf G,\underline{\operatorname{Hom}}(\mathbf G,\mathrm P))=0,\qquad
H^n(\mathbf G,\underline{\operatorname{Hom}}(\mathbf G,\mathrm P))=0\quad(n>0).
\]
\end{SGAresult}
\begin{proof}
只须证明 $\underline C^*(\mathbf G,\underline{\operatorname{Hom}}(\mathbf G,\mathrm P))$ 与 $C^*(\mathbf G,\underline{\operatorname{Hom}}(\mathbf G,\mathrm P))$ 在正\Term{degree}{次数}\Term{homotopically-trivial}{同伦平凡}。甚至只须对后者证明，因为前者的相应结论由\Term{base-change}{基变换}即得。\footnote{编者注（50）：以下对原文作了展开。} 对每个 $n\geq0$，如下定义\Term{morphism}{态射}
\[
\sigma:C^{n+1}(\mathbf G,\underline{\operatorname{Hom}}(\mathbf G,\mathrm P))\longrightarrow C^n(\mathbf G,\underline{\operatorname{Hom}}(\mathbf G,\mathrm P)).
\]
设 $f\in C^{n+1}(\mathbf G,\underline{\operatorname{Hom}}(\mathbf G,\mathrm P))$；对每个 $S\in\operatorname{Ob}(C)$ 及 $g_1,\ldots,g_n\in\mathbf G(S)$，$\sigma(f)(g_1,\ldots,g_n)$ 是如下定义的 $\operatorname{Hom}_S(\mathbf G_S,\mathrm P_S)$ 的元素：对每个 $S'\to S$ 与 $x\in\mathbf G(S')$，
\[
\sigma(f)(g_1,\ldots,g_n)(x)=f(x,g_1,\ldots,g_n)(e)\in\mathrm P(S'),
\]
其中 $e$ 表示 $\mathbf G(S')$ 的\Term{identity-element}{单位元}。于是，$\sigma$ 是正\Term{degree}{次数}上的\Term{homotopy-operator}{同伦算子}。事实上，对每个 $g_1,\ldots,g_{n+1}\in\mathbf G(S)$ 和 $x\in\mathbf G(S')$，一方面有
\[
\begin{aligned}
\partial\sigma(f)(g_1,\ldots,g_{n+1})(x)
={}&f(xg_1,g_2,\ldots,g_{n+1})(e)\\
&+\sum_{i=1}^n(-1)^if(x,g_1,\ldots,g_ig_{i+1},\ldots,g_{n+1})(e)\\
&+(-1)^{n+1}f(x,g_1,\ldots,g_n)(e),
\end{aligned}
\]
另一方面有
\[
\begin{aligned}
\sigma(\partial f)(g_1,\ldots,g_{n+1})(x)
={}&(xf(g_1,g_2,\ldots,g_{n+1}))(e)-f(xg_1,g_2,\ldots,g_{n+1})(e)\\
&+\sum_{i=1}^n(-1)^{i+1}f(x,g_1,\ldots,g_ig_{i+1},\ldots,g_{n+1})\\
&+(-1)^{n+2}f(x,g_1,\ldots,g_n)(e),
\end{aligned}
\]
故
\[
(\partial\sigma(f)+\sigma(\partial f))(g_1,\ldots,g_{n+1})(x)=f(g_1,\ldots,g_{n+1})(x),
\]
即对每个 $n\geq0$，$\partial\sigma+\sigma\partial$ 是 $C^{n+1}(\mathbf G,\underline{\operatorname{Hom}}(\mathbf G,\mathrm P))$ 的\Term{identity-map}{恒等映射}。
\end{proof}

\SourcePage{41}
\begin{SGAresult}{注记}{5.2.3}
\footnote{编者注（51）：此注记为增补。} 假设“$C$ 是\Term{small-category}{小范畴}”仅用于保证\Term{derived-functor}{导出函子} $R^n\underline H^0$ 与 $R^nH^0$ 存在。在任何情况下，上述论证都表明，\Term{functor}{函子} $\underline H^n(\mathbf G,\ )$（分别为 $H^n(\mathbf G,\ )$）构成正\Term{degree}{次数}\Term{effaceable}{可抹除}的\Term{cohomological-functor}{上同调函子}，因而是 [Gr57，2.2] 意义下\Term{left-exact}{左正合}\Term{functor}{函子} $\underline H^0(\mathbf G,\ )$（分别为 $H^0(\mathbf G,\ )$）的（右）\Term{satellite-functor}{卫星函子}；若 $(\mathbf G\text{-}\mathbf O\text{-}\mathrm{Mod.})$ 有足够多的\Term{injective-object}{内射对象}（当 $C$ 为\Term{small-category}{小范畴}时正是如此），它们就与\Term{derived-functor}{导出函子}一致（同处，2.3）。
\end{SGAresult}

\textbf{5.3. $G$-$\mathcal O_S$-模的\Term{cohomology}{上同调}。} 设 $S$ 是\Term{scheme}{概形}，$G$ 是 $S$-\Term{group}{群}，$\mathcal F$ 是\Term{quasi-coherent}{拟凝聚} $G$-$\mathcal O_S$-\Term{module}{模}。定义 $G$ 取值于 $\mathcal F$ 的\Term{cohomology-group}{上同调群}为
\[
H^n(G,\mathcal F)=H^n(h_G,\mathbf W(\mathcal F)).
\]
（关于记号，参见 4.6。）
假设 $G$ 在 $S$ 上\Term{affine}{仿射}。由命题 4.6.4，这个\Term{cohomology}{上同调}按下述方式计算：$H^n(G,\mathcal F)$ 是\Term{complex}{复形} $C^*(G,\mathcal F)$ 的第 $n$ 个\Term{homology-group}{同调群}，其第 $n$ 项为
\[
C^n(G,\mathcal F)=\Gamma\bigl(S,\mathcal F\otimes\underbrace{\mathcal A(G)\otimes\cdots\otimes\mathcal A(G)}_{n\text{ 次}}\bigr).
\]
若 $f$（分别为 $a_i$）是 $\mathcal F$（分别为 $\mathcal A(G)$）在 $S$ 的一个\Term{open-subset}{开集}上的\Term{section}{截面}，则
\[
\begin{aligned}
\partial(f\otimes a_1\otimes\cdots\otimes a_n)
={}&\mu_{\mathcal F}(f)\otimes a_1\otimes a_2\otimes\cdots\otimes a_n\\
&+\sum_{i=1}^n(-1)^if\otimes a_1\otimes\cdots\otimes\Delta a_i\otimes\cdots\otimes a_n\\
&+(-1)^{n+1}f\otimes a_1\otimes a_2\cdots\otimes a_n\otimes1,
\end{aligned}
\]
其中 $\Delta:\mathcal A(G)\to\mathcal A(G)\otimes\mathcal A(G)$ 与 $\mu_{\mathcal F}:\mathcal F\to\mathcal F\otimes\mathcal A(G)$ 描述 $\mathcal A(G)$ 的\Term{coalgebra}{余代数}结构及 $\mathcal F$ 的\Term{comodule}{余模}结构。顺便指出，$G$ 取值于 $\mathcal F$ 的\Term{cohomology}{上同调}因而只依赖于 $\mathcal F$ 的\Term{comodule}{余模}结构，特别地，只依赖于 $G$ 的 $S$-\Term{monoid}{幺半群}结构。
特别地，有
\[
H^0(G,\mathcal F)=\Gamma(S,\mathcal F^G),
\]
其中\Term{invariants-sheaf}{不变量层} $\mathcal F^G$ 定义为这样的\Term{sheaf}{层}：它在 $S$ 的\Term{open-subset}{开集} $U$ 上的\Term{section}{截面}，是 $\mathcal F$ 在 $U$ 上的那些\Term{section}{截面}，其在 $U$ 上的每个 $S'$ 中的\Term{inverse-image}{逆像}均在 $G(S')$ 下\Term{invariant}{不变}（参见 4.7.1.2）。

\begin{SGAresult}{\Term{theorem}{定理}}{5.3.1}
设 $S$ 为\Term{affine}{仿射}\Term{scheme}{概形}，$G$ 为在 $S$ 上\Term{affine}{仿射}且\Term{flat}{平坦}的 $S$-\Term{group}{群}。\Term{functor}{函子} $H^n(G,\ )$ 是\Term{quasi-coherent}{拟凝聚} $G$-$\mathcal O_S$-\Term{module}{模}\Term{category}{范畴}上 $H^0(G,\ )$ 的\Term{derived-functor}{导出函子}。
\end{SGAresult}
\SourcePage{42}
\begin{proof}
\footnote{编者注（52）：修改了原文，以表明\Term{category}{范畴} $(G\text{-}\mathcal O_S\text{-}\mathrm{Mod.q.c.})$ 是\Term{abelian-category}{阿贝尔范畴}，并有足够多的\Term{injective-object}{内射对象}。可比较 [Ja03]，第一部分，3.3--3.4、3.9、4.2 及 4.14--4.16（须注意，其中“$k$-\Term{group-scheme}{群概形}”指“\Term{affine}{仿射} $k$-\Term{group-scheme}{群概形}”，参见该处，2.1）。} 由于 $G$ 在 $S$ 上\Term{affine}{仿射}且\Term{flat}{平坦}，由 4.7.2.1，\Term{category}{范畴} $(G\text{-}\mathcal O_S\text{-}\mathrm{Mod.q.c.})$ \Term{equivalence}{等价}于在 $\mathcal O_S$ 上\Term{quasi-coherent}{拟凝聚}的 $\mathcal A(G)$-\Term{comodule}{余模}\Term{category}{范畴} $(\mathcal A(G)\text{-}\mathrm{Comod.q.c.})$，因而是\Term{abelian-category}{阿贝尔范畴}。另一方面，由于 $\mathcal A(G)$ 是\Term{flat}{平坦} $\mathcal O_S$-\Term{module}{模}，每个\Term{functor}{函子} $\mathcal F\mapsto\mathcal F\otimes_{\mathcal O_S}\mathcal A(G)^{\otimes n}$ 都\Term{exact}{正合}；又因 $S$ \Term{affine}{仿射}，便得到 $C^*(G,\ )$ 是 $(G\text{-}\mathcal O_S\text{-}\mathrm{Mod.q.c.})$ 上的\Term{exact}{正合}\Term{functor}{函子}。
用 $\Delta$（分别为 $\eta$）表示 $\mathcal A(G)$ 的\Term{comultiplication}{余乘法}（分别为\Term{augmentation}{增广}）。对每个\Term{quasi-coherent}{拟凝聚} $\mathcal O_S$-\Term{module}{模} $\mathcal P$，记 $\operatorname{Ind}(\mathcal P)=\mathcal P\otimes_{\mathcal O_S}\mathcal A(G)$，并赋予由
\[
\mathrm{id}_{\mathcal P}\otimes\Delta:\mathcal P\otimes_{\mathcal O_S}\mathcal A(G)\longrightarrow\mathcal P\otimes_{\mathcal O_S}\mathcal A(G)\otimes_{\mathcal O_S}\mathcal A(G)
\]
定义的 $\mathcal A(G)$-\Term{comodule}{余模}结构；这定义了一个\Term{functor}{函子} $\operatorname{Ind}:(\mathcal O_S\text{-}\mathrm{Mod.q.c.})\to(G\text{-}\mathcal O_S\text{-}\mathrm{Mod.q.c.})$。
由 4.6.4.1、5.2.0 与 5.2.0.1，有 $G$-$\mathbf O_S$-\Term{module}{模}\Term{isomorphism}{同构}
\[
\mathbf W(\operatorname{Ind}(\mathcal P))\simeq\mathrm E(\mathbf W(\mathcal P))=\underline{\operatorname{Hom}}(G,\mathbf W(\mathcal P)).\tag{*}
\]
在此等同下，\Term{morphism}{态射} $\varepsilon:\mathrm E(\mathbf W(\mathcal P))\to\mathbf W(\mathcal P)$ 对应于 $\mathcal O_S$-\Term{module-morphism}{模态射} $\mathrm{id}_{\mathcal P}\otimes\eta:\operatorname{Ind}(\mathcal P)\to\mathcal P$。
我们已经使用过\Term{functor}{函子} $\mathbf W:(\mathcal O_S\text{-}\mathrm{Mod.})\to(\mathbf O_S\text{-}\mathrm{Mod.})$ 是\Term{fully-faithful}{全忠实}的；按定义 4.7.1，它在 $(G\text{-}\mathcal O_S\text{-}\mathrm{Mod.})$ 上的\Term{restriction}{限制}也如此，即若 $\mathcal M$、$\mathcal M'$ 为 $G$-$\mathcal O_S$-\Term{module}{模}，就有\Term{functorial}{函子性}\Term{isomorphism}{同构}
\[
\operatorname{Hom}_{G\text{-}\mathcal O_S\text{-}\mathrm{Mod.}}(\mathcal M,\mathcal M')\simeq\operatorname{Hom}_{G\text{-}\mathbf O_S\text{-}\mathrm{Mod.}}(\mathbf W(\mathcal M),\mathbf W(\mathcal M')).
\]
因此，由\Term{lemma}{引理} 5.2.0.2 得到下列推论。
\end{proof}
\begin{SGAresult}{推论}{5.3.1.1}
（i）\Term{functor}{函子} $\operatorname{Ind}$ 是\Term{forgetful-functor}{遗忘函子} $(G\text{-}\mathcal O_S\text{-}\mathrm{Mod.q.c.})\to(\mathcal O_S\text{-}\mathrm{Mod.q.c.})$ 的\Term{right-adjoint}{右伴随}。更确切地说，对 $(G\text{-}\mathcal O_S\text{-}\mathrm{Mod.q.c.})$ 的每个\Term{object}{对象} $\mathcal M$，映射 $\mathrm{id}_{\mathcal P}\otimes\eta:\operatorname{Ind}(\mathcal P)\to\mathcal P$ 诱导\Term{bijection}{双射}
\[
\operatorname{Hom}_{G\text{-}\mathcal O_S\text{-}\mathrm{Mod.}}(\mathcal M,\operatorname{Ind}(\mathcal P))\xrightarrow\sim\operatorname{Hom}_{\mathcal O_S}(\mathcal M,\mathcal P).
\]
（ii）因此，若 $\mathcal I$ 为 $(\mathcal O_S\text{-}\mathrm{Mod.q.c.})$ 的\Term{injective-object}{内射对象}，则 $\operatorname{Ind}(\mathcal I)$ 为 $(G\text{-}\mathcal O_S\text{-}\mathrm{Mod.q.c.})$ 的\Term{injective-object}{内射对象}。
\end{SGAresult}
\begin{proof}[\Term{theorem}{定理} 5.3.1 的证明（续）]
设 $\mathcal F$ 为 $G$-$\mathcal O_S$-\Term{module}{模}，$\mu_{\mathcal F}:\mathcal F\to\operatorname{Ind}(\mathcal F)$ 为定义其 $\mathcal A(G)$-\Term{comodule}{余模}结构的映射。由 5.2.0.3（或者由 4.7.2 的公理（CM 1）与（CM 2）），$\mu_{\mathcal F}$ 是 $G$-$\mathcal O_S$-\Term{module-morphism}{模态射}，且 $(\mathrm{id}_{\mathcal F}\otimes\eta)\circ\mu_{\mathcal F}=\mathrm{id}_{\mathcal F}$；因此，作为 $\mathcal O_S$-\Term{module}{模}，$\mathcal F$ 是 $\operatorname{Ind}(\mathcal F)$ 的\Term{direct-summand}{直和因子}；特别地，$\mu_{\mathcal F}$ 是\Term{monomorphism}{单态射}。由（*）与 5.2.2，
\[
H^n(G,\mathbf W(\operatorname{Ind}(\mathcal F)))\simeq H^n(G,\underline{\operatorname{Hom}}_S(G,\mathbf W(\mathcal F)))=0\quad(n>0),
\]
故 $H^n(G,\ )$ 对 $n>0$ \Term{effaceable}{可抹除}。
最后，因 $S$ \Term{affine}{仿射}，$(\mathcal O_S\text{-}\mathrm{Mod.q.c.})$ 有足够多的\Term{injective-object}{内射对象}。取 $\mathcal O_S$-\Term{module}{模}的\Term{monomorphism}{单态射} $\mathcal F\hookrightarrow\mathcal I$，其中 $\mathcal I$ 为 $(\mathcal O_S\text{-}\mathrm{Mod.q.c.})$ 的\Term{injective-object}{内射对象}；由于 $\mathcal A(G)$ 在 $\mathcal O_S$ 上\Term{flat}{平坦}，$\operatorname{Ind}(\mathcal F)$ 是 $\operatorname{Ind}(\mathcal I)$ 的子 $G$-$\mathcal O_S$-\Term{module}{模}，由此得到：
\end{proof}
\SourcePage{43}
\begin{SGAresult}{推论}{5.3.1.2}
\Term{abelian-category}{阿贝尔范畴} $(G\text{-}\mathcal O_S\text{-}\mathrm{Mod.q.c.})$ 有足够多的\Term{injective-object}{内射对象}。
\end{SGAresult}
结合 [Gr57，2.2.1 与 2.3]（在 5.2.1 的证明中已使用），这就完成了\Term{theorem}{定理} 5.3.1 的证明。

\begin{SGAresult}{注记}{5.3.1.3}
也可通过以下计算证明 5.3.1.1。对每个 $G$-$\mathcal O_S$-\Term{module-morphism}{模态射} $\varphi:\mathcal M\to\mathcal P\otimes_{\mathcal O_S}\mathcal A(G)$，对应 $\mathcal O_S$-\Term{morphism}{态射} $(\mathrm{id}_{\mathcal P}\otimes\eta)\circ\varphi:\mathcal M\to\mathcal P$。反之，对每个 $\mathcal O_S$-\Term{morphism}{态射} $f:\mathcal M\to\mathcal P$，对应 $G$-$\mathcal O_S$-\Term{module-morphism}{模态射} $(f\otimes\mathrm{id}_{\mathcal A(G)})\circ\mu_{\mathcal M}:\mathcal M\to\mathcal P\otimes_{\mathcal O_S}\mathcal A(G)$。立即看出
\[
(\mathrm{id}_{\mathcal P}\otimes\eta)\circ(f\otimes\mathrm{id}_{\mathcal A(G)})\circ\mu_{\mathcal M}=(f\otimes\mathrm{id}_{\mathcal O_S})\circ(\mathrm{id}_{\mathcal P}\otimes\eta)\circ\mu_{\mathcal M}=f.
\]
另一方面，对每个 $\varphi$，下图交换：
\[
\begin{tikzcd}
\mathcal M\arrow[r,"{\varphi}"]\arrow[d,"{\mu_{\mathcal M}}"']&\mathcal P\otimes_{\mathcal O_S}\mathcal A(G)\arrow[d,"{\mathrm{id}_{\mathcal P}\otimes\Delta}"]\\
\mathcal M\otimes_{\mathcal O_S}\mathcal A(G)\arrow[r,"{\varphi\otimes\mathrm{id}_{\mathcal A(G)}}"']&\mathcal P\otimes_{\mathcal O_S}\mathcal A(G)\otimes_{\mathcal O_S}\mathcal A(G).
\end{tikzcd}
\]
由此
\[
\begin{aligned}
\bigl(((\mathrm{id}_{\mathcal P}\otimes\eta)\circ\varphi)\otimes\mathrm{id}_{\mathcal A(G)}\bigr)\circ\mu_{\mathcal M}
&=(\mathrm{id}_{\mathcal P}\otimes\eta\otimes\mathrm{id}_{\mathcal A(G)})\circ(\varphi\otimes\mathrm{id}_{\mathcal A(G)})\circ\mu_{\mathcal M}\\
&=(\mathrm{id}_{\mathcal P}\otimes\eta\otimes\mathrm{id}_{\mathcal A(G)})\circ(\mathrm{id}_{\mathcal P}\otimes\Delta)\circ\varphi=\varphi.
\end{aligned}
\]
这证明了 5.3.1.1（i）（而（ii）由它\Term{pushout}{推出}）。
\end{SGAresult}
设 $\mathcal F$ 为 $G$-$\mathcal O_S$-\Term{module}{模}；上文已见，4.7.2 的公理（CM 2）表明，作为 $\mathcal O_S$-\Term{module}{模}，$\mathcal F$ 是 $\mathrm E(\mathcal F)$ 的\Term{direct-summand}{直和因子}。由此得出：
\begin{SGAresult}{命题}{5.3.2}
设 $S$ 为\Term{affine}{仿射}\Term{scheme}{概形}，$G$ 为\Term{affine}{仿射}且\Term{flat}{平坦}的 $S$-\Term{group}{群}；假设\Term{quasi-coherent}{拟凝聚} $G$-$\mathcal O_S$-\Term{module}{模}的每个\Term{exact-sequence}{正合列} $0\to\mathcal F_1\to\mathcal F_2\to\mathcal F_3\to0$，只要作为 $\mathcal O_S$-\Term{module}{模}序列\Term{split}{分裂}，便也作为 $G$-$\mathcal O_S$-\Term{module}{模}序列\Term{split}{分裂}。
则对 $n>0$，\Term{functor}{函子} $H^n(G,\ )$ 都为零（\Term{equivalence}{等价}地，\Term{functor}{函子} $H^0(G,\ )$ \Term{exact}{正合}）。
\end{SGAresult}
\begin{proof}
依假设，$G$-$\mathcal O_S$-\Term{module}{模}序列
\[
0\longrightarrow\mathcal F\longrightarrow\mathrm E(\mathcal F)\longrightarrow\mathrm E(\mathcal F)/\mathcal F\longrightarrow0
\]
\Term{split}{分裂}；故 $\mathcal F$ 是 $\mathrm E(\mathcal F)$ 的 $G$-$\mathcal O_S$-\Term{module}{模}\Term{direct-summand}{直和因子}，而后者的\Term{cohomology}{上同调}为零。
\end{proof}
由此与命题 4.7.4 立即得到：
\begin{SGAresult}{\Term{theorem}{定理}}{5.3.3}
设 $S$ 为\Term{affine}{仿射}\Term{scheme}{概形}，$G$ 为\Term{diagonalizable}{可对角化} $S$-\Term{group}{群}。对每个\Term{quasi-coherent}{拟凝聚} $G$-$\mathcal O_S$-\Term{module}{模} $\mathcal F$，有 $H^n(G,\mathcal F)=0$，$n>0$。
\end{SGAresult}

\SourcePage{44}
\textbf{注记。} 当 $G$ 在 $S$ 上不必\Term{flat}{平坦}时，命题 5.3.2 仍成立；这时证明须用\Term{relative-cohomology}{相对上同调}。\footnote{编者注（53）：编者未尝试展开此注记。}

\SGAsection{6}{$G$-等变对象与模}
\footnote{编者注（54）：本节为增补。} 设 $C$ 为一个有\Term{final-object}{终对象}且存在\Term{fiber-product}{纤维积}的\Term{category}{范畴}。设 $\mathbf G$ 为 $\widehat C$-\Term{group}{群}，$\pi:\mathrm M\to\mathrm X$ 为 $\widehat C$ 的\Term{morphism}{态射}，且 $\lambda=\lambda_{\mathrm X}:\mathbf G\times\mathrm X\to\mathrm X$ 为 $\mathbf G$ 在 $\mathrm X$ 上的\Term{action}{作用}。下文中，用 $\mathrm Y\times_f\mathrm M$ 表示 $\pi:\mathrm M\to\mathrm X$ 与一个 $\mathrm X$-\Term{functor}{函子} $f:\mathrm Y\to\mathrm X$ 的\Term{fiber-product}{纤维积}。
对每个 $U\in\operatorname{Ob}(C)$ 与 $x\in\mathrm X(U)$，记 $\mathrm M_x=U\times_x\mathrm M$，即对每个 $\varphi:U'\to U$，有
\[
\mathrm M_x(U')=\{m\in\mathrm M(U')\mid\pi(m)=x_{U'}=\varphi^*(x)\}.
\]
最后，若 $g\in\mathbf G(U)$，也用 $gx$ 表示 $\mathrm X(U)$ 的元素 $\lambda(g,x)$。
\begin{SGAresult}{定义}{6.1}
a）若在 $\mathrm M$ 上给定一个\Term{lifting}{提升} $\lambda$ 的 $\mathbf G$-\Term{action}{作用} $\Lambda:\mathbf G\times\mathrm M\to\mathrm M$，即使以下方块交换的\Term{action}{作用}，就称 $\mathrm M$ 为 $\mathbf G$-\Term{equivariant}{等变} $\mathrm X$-\Term{object}{对象}：
\[
\begin{tikzcd}
\mathbf G\times\mathrm M\arrow[r,"{\Lambda}"]\arrow[d,"{\mathrm{id}_{\mathbf G}\times\pi}"']&\mathrm M\arrow[d,"{\pi}"]\\
\mathbf G\times\mathrm X\arrow[r,"{\lambda}"']&\mathrm X.
\end{tikzcd}
\]
这\Term{equivalence}{等价}于对每个\Term{morphism}{态射} $(g,x):U\to\mathbf G\times\mathrm X$ 给定映射
\[
\Lambda_x^U(g):\mathrm M_x(U)\longrightarrow\mathrm M_{gx}(U),\qquad m\longmapsto g\cdot m,
\]
满足 $1\cdot m=m$ 与 $g\cdot(h\cdot m)=(gh)\cdot m$，并关于 $(\mathbf G\times\mathrm X)$-\Term{object}{对象} $U$ \Term{functorial}{函子性}。这又\Term{equivalence}{等价}于给定 $U$-\Term{object}{对象}\Term{morphism}{态射}
\[
\Lambda_x(g):\mathrm M_x\longrightarrow\mathrm M_{gx},
\]
满足 $\Lambda_x(1)=\mathrm{id}$ 与 $\Lambda_{hx}(g)\circ\Lambda_x(h)=\Lambda_x(gh)$。

b）现设 $\mathbf O$ 为 $\widehat C$-\Term{ring}{环}，$\mathbf O_{\mathrm X}=\mathbf O\times\mathrm X$。在（a）的条件下，若 $\mathrm M$ 是 $\mathbf O_{\mathrm X}$-\Term{module}{模}（参见定义 4.3.3.1，它对任意\Term{category}{范畴} $C$ 上的\Term{ring}{环}\Term{functor}{函子}均适用），且\Term{action}{作用} $\Lambda$ 与 $\mathrm M$ 的 $\mathbf O_{\mathrm X}$-\Term{module}{模}结构相容，即对每个 $(g,x)\in\mathbf G(U)\times\mathrm X(U)$（$U\in\operatorname{Ob}(C)$），映射 $\Lambda_x(g):\mathrm M_x\to\mathrm M_{gx}$、$m\mapsto g\cdot m$ 是 $\mathbf O_U$-\Term{module-morphism}{模态射}，就称 $\mathrm M$ 为 $\mathbf G$-\Term{equivariant}{等变} $\mathbf O_{\mathrm X}$-\Term{module}{模}。
\end{SGAresult}
\begin{SGAresult}{注记}{6.2}
（i）在上述（a）中，条件 $\Lambda_x(1)=\mathrm{id}$ 与 $\Lambda_{hx}(g)\circ\Lambda_x(h)=\Lambda_x(gh)$ 显然意味着每个 $\Lambda_x(g)$ 都是\Term{isomorphism}{同构}，其逆为 $\Lambda_{gx}(g^{-1})$。反之，若假设每个 $\Lambda_x(g)$ 都是\Term{isomorphism}{同构}，则在条件 $\Lambda_{hx}(g)\circ\Lambda_x(h)=\Lambda_x(gh)$ 中取 $h=1$，就得到 $\Lambda_x(1)=\mathrm{id}$。
\SourcePage{45}
（ii）设 $\mathrm M$ 为 $\mathbf O_{\mathrm X}$-\Term{module}{模}。首先可见，给定一个使 6.1 的图交换且使每个\Term{morphism}{态射} $\Lambda_x(g):\mathrm M_x\to\mathrm M_{gx}$、$m\mapsto g\cdot m$ 都为 $\mathbf O_U$-\Term{module}{模}\Term{isomorphism}{同构}的\Term{morphism}{态射} $\Lambda:\mathbf G\times\mathrm M\to\mathrm M$，\Term{equivalence}{等价}于给定一个 $\mathbf O_{\mathbf G\times\mathrm X}$-\Term{module}{模}\Term{isomorphism}{同构}
\[
\begin{aligned}
\theta:\mathbf G\times\mathrm M=(\mathbf G\times\mathrm X)\times_{\operatorname{pr}_{\mathrm X}}\mathrm M&\xrightarrow\sim(\mathbf G\times\mathrm X)\times_\lambda\mathrm M,\\
(g,x,m)&\longmapsto(g,x,g\cdot m).
\end{aligned}
\]
因已假设每个 $\Lambda_x(g)$ 为\Term{isomorphism}{同构}，在等式 $\Lambda_{hx}(g)\circ\Lambda_x(h)=\Lambda_x(gh)$ 中取 $h=1$ 就能\Term{pushout}{推出} $\Lambda_x(1)=\mathrm{id}$。所以 $\Lambda$ 是 $\mathbf G$ 在 $\mathrm M$ 上的\Term{action}{作用}（即 $g\cdot(h\cdot m)=(gh)\cdot m$），当且仅当以下 $(\mathbf G\times\mathbf G\times\mathrm X)$-\Term{isomorphism}{同构}的图交换（其中 $\mu$ 表示 $\mathbf G$ 的\Term{multiplication}{乘法}，$f^*(\theta)$ 表示由 $\theta$ 作\Term{base-change}{基变换} $f:\mathbf G\times\mathbf G\times\mathrm X\to\mathbf G\times\mathrm X$ 所得的\Term{isomorphism}{同构}）：
\[
\begin{tikzcd}[column sep=small,row sep=large,font=\small]
(\mathbf G\times\mathbf G\times\mathrm X)\times_{\operatorname{pr}_{\mathrm X}\circ\operatorname{pr}_{2,3}}\mathrm M\arrow[r,"{\operatorname{pr}_{2,3}^*(\theta)}","{\sim}"']\arrow[d,equal] & (\mathbf G\times\mathbf G\times\mathrm X)\times_{\lambda\circ\operatorname{pr}_{2,3}}\mathrm M\arrow[d,equal]\\
(\mathbf G\times\mathbf G\times\mathrm X)\times_{\operatorname{pr}_{\mathrm X}\circ(\mu\times\mathrm{id}_{\mathrm X})}\mathrm M\arrow[d,"{(\mu\times\mathrm{id}_{\mathrm X})^*(\theta)}"',"{\sim}"] & (\mathbf G\times\mathbf G\times\mathrm X)\times_{\operatorname{pr}_{\mathrm X}\circ(\mathrm{id}_{\mathbf G}\times\lambda)}\mathrm M\arrow[d,"{(\mathrm{id}_{\mathbf G}\times\lambda)^*(\theta)}","{\sim}"']\\
(\mathbf G\times\mathbf G\times\mathrm X)\times_{\lambda\circ(\mu\times\mathrm{id}_{\mathrm X})}\mathrm M\arrow[r,equal] & (\mathbf G\times\mathbf G\times\mathrm X)\times_{\lambda\circ(\mathrm{id}_{\mathbf G}\times\lambda)}\mathrm M.
\end{tikzcd}
\]
因而可见，在 $\mathrm M$ 上给定 $\mathbf G$-\Term{equivariant}{等变} $\mathbf O_{\mathrm X}$-\Term{module}{模}结构，\Term{equivalence}{等价}于给定一个如上的 $\mathbf O_{\mathbf G\times\mathrm X}$-\Term{module}{模}\Term{isomorphism}{同构} $\theta$，使上图交换。

（iii）上述一切可推广到 $\mathbf G$ 仅为 $\widehat C$-\Term{monoid}{幺半群}的情形：此时，给定一个\Term{lifting}{提升} $\lambda$、且使每个 $\Lambda_x(g):\mathrm M_x\to\mathrm M_{gx}$ 都为 $\mathbf O_U$-\Term{module-morphism}{模态射}的\Term{action}{作用} $\Lambda:\mathbf G\times\mathrm M\to\mathrm M$，\Term{equivalence}{等价}于给定一个如（ii）的 $\mathbf O_{\mathbf G\times\mathrm X}$-\Term{module-morphism}{模态射} $\theta$，使上图（去掉\Term{arrow}{箭头}下的 $\sim$ \Term{signature}{符号}）交换，且满足 $p_{\mathrm M}\circ\theta\circ(\varepsilon_{\mathbf G}\times\mathrm{id}_{\mathrm M})=\mathrm{id}_{\mathrm M}$，其中 $\varepsilon_{\mathbf G}$ 表示 $\mathbf G$ 的\Term{unit-section}{单位截面}，$p_{\mathrm M}$ 表示到 $\mathrm M$ 的\Term{projection}{投影}。
\end{SGAresult}

\textbf{6.3. $\mathbf G$-\Term{equivariant}{等变}\Term{morphism}{态射}。} 设 $\mathrm Y$ 为 $\widehat C$ 的另一个\Term{object}{对象}，带有 $\mathbf G$ 的\Term{action}{作用} $\lambda_{\mathrm Y}:\mathbf G\times\mathrm Y\to\mathrm Y$；设 $\mathrm N$ 为另一个 $\mathbf G$-\Term{equivariant}{等变} $\mathbf O_{\mathrm X}$-\Term{module}{模}。若 $\widehat C$-\Term{morphism}{态射} $f:\mathrm Y\to\mathrm X$（分别为 $\mathbf O_{\mathrm X}$-\Term{module-morphism}{模态射} $\varphi:\mathrm M\to\mathrm N$）与 $\mathbf G$ 的\Term{action}{作用}交换，即在\Term{set}{集合}意义下满足 $f(g\cdot y)=g\cdot f(y)$（分别为 $\varphi(g\cdot m)=g\cdot\varphi(m)$），就称其为 $\mathbf G$-\Term{equivariant}{等变}的；这\Term{equivalence}{等价}于 $f\circ\lambda_{\mathrm Y}=\lambda_{\mathrm X}\circ(\mathrm{id}_{\mathbf G}\times f)$（分别为 $\varphi\circ\Lambda_{\mathrm M}=\Lambda_{\mathrm N}\circ(\mathrm{id}_{\mathbf G}\times\varphi)$）。立即得到下列\Term{lemma}{引理}：

\SourcePage{46}
\begin{SGAresult}{\Term{lemma}{引理}}{6.3.1}
设 $f:\mathrm Y\to\mathrm X$ 为 $\mathbf G$-\Term{equivariant}{等变}\Term{morphism}{态射}，$\mathrm M$ 为 $\mathbf G$-\Term{equivariant}{等变} $\mathbf O_{\mathrm X}$-\Term{module}{模}。则\Term{inverse-image}{逆像} $f^*(\mathrm M)=\mathrm Y\times_f\mathrm M$ 是 $\mathbf G$-\Term{equivariant}{等变} $\mathbf O_{\mathrm Y}$-\Term{module}{模}。
\end{SGAresult}
另一方面，$\mathbf G$ \Term{action}{作用}于 $\underline{\operatorname{Hom}}_{\widehat C}(\mathrm Y,\mathrm X)$。事实上，设 $T\in\operatorname{Ob}(C)$，$g\in\mathbf G(T)$，$\varphi$ 为 $T$-\Term{morphism}{态射} $\mathrm Y_T\to\mathrm X_T$；则 $g$ 定义 $\mathrm Y_T$ 和 $\mathrm X_T$ 的\Term{automorphism}{自同构} $\lambda_{\mathrm Y}(g)$ 与 $\lambda_{\mathrm X}(g)$，并用 $g\cdot\varphi$（也写作 $g\varphi g^{-1}$）表示\Term{morphism}{态射} $\lambda_{\mathrm X}(g)\circ\varphi\circ\lambda_{\mathrm Y}(g^{-1})$。这定义了 $\mathbf G(T)$ 在 $\operatorname{Hom}_T(\mathrm Y_T,\mathrm X_T)$ 上关于 $T$ \Term{functorial}{函子性}的\Term{action}{作用}。
\begin{SGAresult}{定义}{6.3.2}
因此，若 $\varphi:\mathrm Y\to\mathrm X$ 为任意 $\widehat C$-\Term{morphism}{态射}，就可以考虑其\Term{stabilizer}{稳定子} $\operatorname{Stab}_{\mathbf G}(\varphi)$（参见 2.3.3.1）：对每个 $T\in\operatorname{Ob}(C)$，$\operatorname{Stab}_{\mathbf G}(\varphi)(T)$ 是 $\mathbf G(T)$ 中由满足 $g\circ\varphi_T=\varphi_T\circ g$ 的 $g$ 所组成的\Term{subgroup}{子群}，即使下图交换的 $g$：
\[
\begin{tikzcd}
\mathrm Y_T\arrow[r,"{\varphi_T}"]\arrow[d,"{g}"']&\mathrm X_T\arrow[d,"{g}"]\\
\mathrm Y_T\arrow[r,"{\varphi_T}"']&\mathrm X_T.
\end{tikzcd}
\]
于是，\Term{morphism}{态射} $\varphi:\mathrm Y\to\mathrm X$ 在 $\operatorname{Stab}_{\mathbf G}(\varphi)$ 下\Term{equivariant}{等变}。
\end{SGAresult}

\textbf{6.4. \Term{global-section}{整体截面}。} 设 $\mathrm M$ 为 $\mathbf G$-\Term{equivariant}{等变} $\mathbf O_{\mathrm X}$-\Term{module}{模}。以 $S_0$ 表示 $C$ 的\Term{final-object}{终对象}，并以 $\prod_{\mathrm X/S_0}\mathrm M$（参见第 II 讲，1.1）表示“$\mathrm M$ 在 $\mathrm X$ 上的\Term{section}{截面}\Term{functor}{函子}”：它向每个 $T\in\operatorname{Ob}(C)$ 赋予
\[
\operatorname{Hom}_{\mathrm X}(\mathrm X_T,\mathrm M)=\operatorname{Hom}_{\mathrm X_T}(\mathrm X_T,\mathrm M_T)=\Gamma(\mathrm M_T/\mathrm X_T).
\]
另一方面回忆，$\mathrm X$ 上的 $\widehat C$-\Term{object}{对象}的每个\Term{morphism}{态射} $g:\mathrm Z\to\mathrm Y$ 都诱导一个\Term{abelian-group}{交换群}\Term{morphism}{态射} $\mathrm M(g):\mathrm M(\mathrm Y)\to\mathrm M(\mathrm Z)$，它与\Term{ring-morphism}{环态射} $g^*:\mathbf O(\mathrm Y)\to\mathbf O(\mathrm Z)$ 相容。特别地，当 $\mathrm Z=\mathrm Y$（于是 $g$ 为 $\mathrm Y$ 的 $\mathrm X$-\Term{endomorphism}{自同态}）时，得到\Term{abelian-group}{交换群}\Term{morphism}{态射}
\[
\mathrm M(g):\mathrm M(\mathrm Y)\longrightarrow\mathrm M(\mathrm Y),
\]
它一般并非 $\mathbf O(\mathrm Y)$-\Term{module-morphism}{模态射}，但对每个 $m\in\mathrm M(\mathrm Y)$ 与 $\alpha\in\mathbf O(\mathrm Y)$，满足
\[
\mathrm M(g)(\alpha\cdot m)=g^*(\alpha)\cdot\mathrm M(g)(m).
\]
下文中简写 $g^*$ 代替 $\mathrm M(g)$。
设 $T\in\operatorname{Ob}(C)$，$\mathrm X_T=\mathrm X\times T$，$\operatorname{pr}_{\mathrm X}$ 为\Term{projection}{投影} $\mathrm X_T\to\mathrm X$。对每个 $\alpha\in\mathbf O(\mathrm X_T)$ 与 $g\in\mathbf G(T)$，令 $g(\alpha)=(g^{-1})^*(\alpha)$：这是 $\mathbf O(\mathrm X_T)$ 中按下述\Term{set}{集合}式公式定义的元素：对每个 $S\in\operatorname{Ob}(C)$ 与 $x\in\mathrm X(S)$、$t\in T(S)$，
\[
g(\alpha)(x,t)=\alpha(g^{-1}x,t).
\]
由此得到 $\mathbf G(T)$ 以\Term{ring}{环}\Term{automorphism}{自同构}在 $\mathbf O(\mathrm X_T)$ 上的（左）\Term{action}{作用}，它关于 $T$ \Term{functorial}{函子性}，且当 $\alpha$ 为 $\mathbf O(T)$ 中某个元素在 $\mathbf O(\mathrm X_T)$ 中的像时，有 $g(\alpha)=\alpha$。
现在令 $\varphi$ 为 $\mathrm X_T$ 的恒等\Term{morphism}{态射}（参见 6.3 及下文对\Term{morphism}{态射} $\varphi:\mathrm Y\to\mathrm X$ 的推广），并根据把 $\mathrm X_T$ 通过 $\operatorname{pr}_{\mathrm X}\circ\lambda(g^{-1})_T$ 或通过 $\operatorname{pr}_{\mathrm X}$ 看成 $\mathrm X$-\Term{object}{对象}，分别用 $\mathrm M(\varphi g^{-1})$ 或 $\mathrm M(\varphi)$ 表示 $\operatorname{Hom}_{\mathrm X}(\mathrm X_T,\mathrm M)$。
\SourcePage{47}
于是 $\lambda(g^{-1})_T$ 是这两个 $\mathrm X$-\Term{object}{对象}之间的 $\mathrm X$-\Term{morphism}{态射}；由上述讨论，得到\Term{abelian-group}{交换群}\Term{morphism}{态射}
\[
(g^{-1})^*:\mathrm M(\varphi)\longrightarrow\mathrm M(\varphi g^{-1}),\qquad m\longmapsto m\circ\lambda(g^{-1})_T,
\]
对每个 $m\in\mathrm M(\varphi)$ 与 $\alpha\in\mathbf O(\mathrm X_T)$，它满足 $(g^{-1})^*(\alpha\cdot m)=(g\alpha)\cdot(g^{-1})^*(m)$。（若 $m$ 是 $\mathrm M_T$ 在 $\mathrm X_T$ 上的\Term{section}{截面}，则 $(g^{-1})^*(m)$ 是由\Term{set}{集合}式公式 $(x,t)\mapsto m(g^{-1}x,t)$ 定义的\Term{section}{截面}。）特别地，$(g^{-1})^*$ 是 $\mathbf O(T)$-\Term{module-morphism}{模态射}。
由 $\mathbf O(\mathrm X_T)$-\Term{module-morphism}{模态射} $\Lambda_x(g)$ 的\Term{functorial}{函子性}，得到\Term{commutative-diagram}{交换图}
\[
\begin{tikzcd}
\mathrm M(\varphi)\arrow[r,"{(g^{-1})^*}"]\arrow[d,"{\Lambda_\varphi(g)}"']&\mathrm M(\varphi g^{-1})\arrow[d,"{\Lambda_{\varphi g^{-1}}(g)}"]\\
\mathrm M(g\varphi)\arrow[r,"{(g^{-1})^*}"']&\mathrm M(g\varphi g^{-1}).
\end{tikzcd}
\]
由于 $\varphi$ 是\Term{identity-map}{恒等映射}，$g\varphi g^{-1}=\varphi$。记 $A(g)=(g^{-1})^*\circ\Lambda_\varphi(g)$，因而得到\Term{abelian-group}{交换群}\Term{morphism}{态射}
\[
A(g):\mathrm M(\varphi)\longrightarrow\mathrm M(\varphi),
\]
它“与 $\mathbf G(T)$ 在 $\mathbf O(\mathrm X_T)$ 上的\Term{action}{作用}相容”，即满足
\[
A(g)(\alpha\cdot m)=(g\alpha)\cdot A(g)(m).
\]
最后，若 $h$ 为 $\mathbf G(T)$ 的另一个元素，由 $\mathrm M$ 及\Term{morphism}{态射} $\Lambda_x(g)$ 的\Term{functorial}{函子性}，有\Term{commutative-diagram}{交换图}
\[
\begin{tikzcd}[column sep=large,row sep=large]
\mathrm M(\varphi)\arrow[d,"{\Lambda_\varphi(g)}"']\arrow[dr,"{\Lambda_\varphi(hg)}"]&&\\
\mathrm M(g\varphi)\arrow[r,"{\Lambda_\varphi(h)}"]\arrow[d,"{(g^{-1})^*}"']&\mathrm M(hg\varphi)\arrow[d,"{(g^{-1})^*}"']\arrow[dr,"{((hg)^{-1})^*}"]&\\
\mathrm M(\varphi)\arrow[r,"{\Lambda_\varphi(h)}"']&\mathrm M(h\varphi)\arrow[r,"{(h^{-1})^*}"']&\mathrm M(\varphi),
\end{tikzcd}
\]
由此 $A(h)\circ A(g)=A(hg)$。因此得到了以下命题。
\begin{SGAresult}{命题}{6.4.1}
对每个 $T$，$\mathbf O(\mathrm X_T)$-\Term{module}{模} $(\prod_{\mathrm X/S_0}\mathrm M)(T)=\Gamma(\mathrm M_T/\mathrm X_T)$ 都带有一个关于 $T$ \Term{functorial}{函子性}的 $\mathbf G(T)$-\Term{action}{作用}，它“与 $\mathbf G(T)$ 在 $\mathbf O(\mathrm X_T)$ 上的\Term{action}{作用}相容”，即满足
\[
A(g)(\alpha\cdot m)=g(\alpha)\cdot A(g)(m).
\]
由于对 $\alpha\in\mathbf O(T)$ 有 $g(\alpha)=\alpha$，特别得到 $\prod_{\mathrm X/S_0}\mathrm M$ 是 $\mathbf G$-$\mathbf O_{S_0}$-\Term{module}{模}。
\end{SGAresult}
更一般地，如 6.3，设 $\mathrm Y$ 为另一个 $\mathbf G$-\Term{equivariant}{等变} $\widehat C$-\Term{object}{对象}，$\varphi:\mathrm Y\to\mathrm X$ 为 $\widehat C$-\Term{morphism}{态射}，$\mathbf H=\operatorname{Stab}_{\mathbf G}(\varphi)$。则\Term{fiber-product}{纤维积} $\mathrm M_\varphi=\mathrm Y\times_\varphi\mathrm M$ 是 $\mathbf H$-\Term{equivariant}{等变} $\mathbf O_{\mathrm Y}$-\Term{module}{模}。因此得到：
\SourcePage{48}
\begin{SGAresult}{推论}{6.4.2}
\Term{functor}{函子} $\prod_{\mathrm Y/S_0}\mathrm M_\varphi$ 是 $\operatorname{Stab}_{\mathbf G}(\varphi)$-$\mathbf O_{S_0}$-\Term{module}{模}。
\end{SGAresult}

\textbf{6.5. $G$-\Term{equivariant}{等变} $\mathcal O_X$-模。} 把上述论述应用于如下情形。设 $S$ 为\Term{scheme}{概形}，$G$ 为\Term{action}{作用}在 $S$-\Term{scheme}{概形} $X$ 上的 $S$-\Term{group-scheme}{群概形}，$\mathcal F$ 为 $\mathcal O_X$-\Term{module}{模}（不必\Term{quasi-coherent}{拟凝聚}）。
\begin{SGAresult}{定义}{6.5.1}
若 $\mathbf O_X$-\Term{module}{模} $\mathrm M=\mathbf W(\mathcal F)$ 是 $G$-\Term{equivariant}{等变}的，就称 $\mathcal F$ 为 $G$-\Term{equivariant}{等变} $\mathcal O_X$-\Term{module}{模}；也就是说，对每个\Term{morphism}{态射} $(g,x):U\to G\times_S X$，给定 $\mathbf O_U$-\Term{module}{模}\Term{isomorphism}{同构}
\[
\Lambda_x(g):\mathbf W(x^*(\mathcal F))\xrightarrow\sim\mathbf W((gx)^*(\mathcal F)),
\]
它们关于 $U$ \Term{functorial}{函子性}，并满足 $\Lambda_{hx}(g)\circ\Lambda_x(h)=\Lambda_x(gh)$。因\Term{functor}{函子} $\mathbf W$ \Term{fully-faithful}{全忠实}（参见 4.6.1.1 与 4.6.2），便得到 $\mathcal O_U$-\Term{module}{模}\Term{isomorphism}{同构}
\[
\Lambda_x(g):x^*(\mathcal F)\xrightarrow\sim(gx)^*(\mathcal F),
\]
这里回忆 $gx$ 表示\Term{morphism}{态射} $\lambda\circ(g,x):U\to X$。特别地，将此应用于 $G\times_S X$ 的恒等\Term{morphism}{态射}，得到 $\mathcal O_{G\times_S X}$-\Term{module}{模}\Term{isomorphism}{同构}
\[
\theta:\operatorname{pr}_X^*(\mathcal F)\xrightarrow\sim\lambda^*(\mathcal F)\tag{$\star$}
\]
使以下 $G\times_S G\times_S X$ 上的\Term{sheaf}{层}\Term{morphism}{态射}图（$\star\star$）交换：
\[
\begin{tikzcd}[column sep=small,row sep=large,font=\small]
\operatorname{pr}_{2,3}^*\circ\operatorname{pr}_X^*(\mathcal F)\arrow[r,"{\operatorname{pr}_{2,3}^*(\theta)}","{\sim}"']\arrow[d,equal]&\operatorname{pr}_{2,3}^*\circ\lambda^*(\mathcal F)\arrow[r,equal]&(\mathrm{id}_G\times\lambda)^*\circ\operatorname{pr}_X^*(\mathcal F)\arrow[d,"{(\mathrm{id}_G\times\lambda)^*(\theta)}","{\sim}"']\\
(\mu\times\mathrm{id}_X)^*\circ\operatorname{pr}_X^*(\mathcal F)\arrow[r,"{(\mu\times\mathrm{id}_X)^*(\theta)}"',"{\sim}"]&(\mu\times\mathrm{id}_X)^*\circ\lambda^*(\mathcal F)\arrow[r,equal]&(\mathrm{id}_G\times\lambda)^*\circ\lambda^*(\mathcal F).
\end{tikzcd}\tag{$\star\star$}
\]
此外，若 $\mathcal E$ 为另一个 $G$-\Term{equivariant}{等变} $\mathcal O_X$-\Term{module}{模}，则当\Term{morphism}{态射} $\mathbf W(\varphi):\mathbf W(\mathcal E)\to\mathbf W(\mathcal F)$ 为 $G$-\Term{equivariant}{等变}时，称 $\mathcal O_X$-\Term{module-morphism}{模态射} $\varphi:\mathcal E\to\mathcal F$ 为 $G$-\Term{equivariant}{等变}。用上述记号，这\Term{equivalence}{等价}于 $\theta_{\mathcal F}\circ\operatorname{pr}_X^*(\varphi)=\lambda^*(\varphi)\circ\theta_{\mathcal E}$。
\end{SGAresult}
\begin{SGAresult}{注记}{6.5.2}
若 $f:Y\to X$ 为 $G$-\Term{equivariant}{等变}\Term{morphism}{态射}，则由 6.3.1，$f^*(\mathcal F)$ 为 $G$-\Term{equivariant}{等变} $\mathcal O_Y$-\Term{module}{模}。
\end{SGAresult}
\begin{SGAresult}{注记}{6.5.3}
再假设 $\mathcal F$ \Term{quasi-coherent}{拟凝聚}。由上述论述，在 $\mathrm M=\mathbf W(\mathcal F)$ 上给定 $G$-\Term{equivariant}{等变} $\mathbf O_X$-\Term{module}{模}结构，\Term{equivalence}{等价}于在 $\mathbf V(\mathcal F)$ 上给定这样的结构，又\Term{equivalence}{等价}于在\Term{vector-fibration}{向量纤维化} $\mathbb V(\mathcal F)$ 上给定一个与 $X$ 上的\Term{action}{作用}相容且在\Term{fibre}{纤维}上“\Term{linear}{线性}”的 $G$-\Term{action}{作用}。
事实上，用 $\varphi$ 表示 $\theta$ 的逆\Term{isomorphism}{同构} $\lambda^*(\mathcal F)\xrightarrow\sim\operatorname{pr}_X^*(\mathcal F)$。对每个\Term{morphism}{态射} $(g,x):U\to G\times_S X$，有一个从 $(gx)^*(\mathcal F)$ 到 $x^*(\mathcal F)$ 的 $\mathcal O_U$-\Term{module}{模}\Term{isomorphism}{同构} $\varphi_x(g)=\Lambda_x(g)^{-1}=\Lambda_{gx}(g^{-1})$，它诱导 $\mathbf O_U$-\Term{module}{模}\Term{isomorphism}{同构}
\[
\mathbf V((gx)^*(\mathcal F))\xrightarrow\sim\mathbf V(x^*(\mathcal F)),
\]
\SourcePage{49}
记为 ${}^t\varphi_x(g)$。于是有
\[
{}^t\varphi_{gx}(h)\circ{}^t\varphi_x(g)={}^t\bigl(\Lambda_{gx}(h)\circ\varphi_x(g)\bigr)^{-1}={}^t\Lambda_x(hg)^{-1}={}^t\varphi_x(hg).
\]
因为对每个 $X$-\Term{scheme}{概形} $f:Y\to X$，都有 $\mathbf V(\mathcal F)\times_fY\simeq\mathbf V(f^*(\mathcal F))$（参见 4.6.2），所以\Term{isomorphism}{同构}
\[
\begin{aligned}
{}^t\varphi:\ G\times_S\mathbf V(\mathcal F)&=(G\times_S X)\times_{\operatorname{pr}_X}\mathbf V(\mathcal F)=\mathbf V(\operatorname{pr}_X^*(\mathcal F))\\
&\xrightarrow\sim\mathbf V(\lambda^*(\mathcal F))=(G\times_S X)\times_\lambda\mathbf V(\mathcal F)
\end{aligned}
\]
在 $\mathbf V(\mathcal F)$ 上给出 $G$-\Term{equivariant}{等变} $\mathbf O_X$-\Term{module}{模}结构。最后，把每个\Term{functor}{函子} $\mathbf V(-)$ 等同于表示它的\Term{vector-fibration}{向量纤维化}，并把 ${}^t\varphi$ 与到 $\mathbb V(\mathcal F)$ 的\Term{projection}{投影}复合，就得到 $G$ 在\Term{vector-fibration}{向量纤维化} $\mathbb V(\mathcal F)$ 上的\Term{action}{作用}，它与 $X$ 上的\Term{action}{作用}相容，且在\Term{fibre}{纤维}上“\Term{linear}{线性}”。
这样便重新得到例如 [GIT，第 1 章，§3] 所给的定义，差别在于该处 Mumford 考虑的是\Term{locally-free}{局部自由}\Term{finite-rank}{有限秩} $\mathcal O_X$-\Term{module}{模} $\mathcal E$，并在 $\mathbf V(\mathcal E)=\mathbf W(\mathcal E^\vee)$ 上定义 $G$-\Term{action}{作用}。事实上，把前面的图（$\star\star$）中 $\theta$ 换成 $\varphi$ 并反转\Term{arrow}{箭头}方向，便恰好得到该处定义 1.6 的图；上述\Term{isomorphism}{同构} ${}^t\varphi$ 也与该处第 31 页的\Term{isomorphism}{同构} $\Phi$ 一致。
\end{SGAresult}
\begin{SGAresult}{注记}{6.5.4}
特别考虑 $X=S$ 且带有 $G$ 的\Term{trivial-action}{平凡作用}的情形。这时，$G$-\Term{equivariant}{等变} $\mathcal O_S$-\Term{module}{模} $\mathcal F$ 与 $G$-$\mathcal O_S$-\Term{module}{模}是同一回事（参见 4.7.1）。此外，若以 $f$ 表示\Term{morphism}{态射} $G\to S$（这里它等于 $\operatorname{pr}_X$，也等于 $\lambda$），则\Term{isomorphism}{同构}
\[
\theta:f^*(\mathcal F)\xrightarrow\sim f^*(\mathcal F)\tag{$\star$}
\]
是下式中的一个元素：
\[
\begin{aligned}
\operatorname{Hom}_{\mathcal O_G}(f^*(\mathcal F),f^*(\mathcal F))
&=\operatorname{Hom}_{\mathbf O_G}(\mathbf W(f^*(\mathcal F)),\mathbf W(f^*(\mathcal F)))\\
&=\underline{\operatorname{End}}_{\mathbf O_S}(\mathbf W(\mathcal F))(G),
\end{aligned}
\]
它就是定义 $G$ 在 $\mathbf W(\mathcal F)$ 上的\Term{action}{作用}的\Term{morphism}{态射} $\rho:G\to\underline{\operatorname{End}}_{\mathbf O_S}(\mathbf W(\mathcal F))$。此外，$\theta$ 通过\Term{adjunction}{伴随}对应于 $\mathcal O_S$-\Term{module-morphism}{模态射} $f_*(\theta)\circ\tau:\mathcal F\to f_*f^*(\mathcal F)$，其中 $\tau:\mathcal F\to f_*f^*(\mathcal F)$ 为\Term{adjunction}{伴随}的“单位”\Term{morphism}{态射}。（这将在 $\mathrm{VI_B}$，11.10.bis 中使用。）
\end{SGAresult}

\textbf{6.6. \Term{functor}{函子} $\prod_{X/S}\mathbf W(\mathcal F)$ 与 $\mathbf W(p_*(\mathcal F))$。} 设 $S$ 为\Term{scheme}{概形}，$G$ 为通过\Term{morphism}{态射} $\lambda:G\times_S X\to X$ 与 $\mu:G\times_S Y\to Y$ \Term{action}{作用}于 $S$-\Term{scheme}{概形} $X$、$Y$ 的 $S$-\Term{group-scheme}{群概形}，且 $p:X\to Y$ 为 $G$-\Term{equivariant}{等变}\Term{morphism}{态射}。假设\Term{morphism}{态射} $p:X\to Y$ 与 $\nu_Y:Y\to S$ 都\Term{quasi-compact}{拟紧}且\Term{quasi-separated}{拟分离}，并且 $\pi:G\to S$ \Term{flat}{平坦}。
于是\Term{projection}{投影} $\operatorname{pr}_Y:G\times_S Y\to Y$ \Term{flat}{平坦}，$\mu$ 也\Term{flat}{平坦}（因为 $\mu$ 是 $\operatorname{pr}_Y$ 与\Term{automorphism}{自同构} $(g,y)\mapsto(g,gy)$ 的复合）。最后，对可变的 $S$-\Term{scheme}{概形} $f:T\to{}$，用 $p_T:X_T\to T$ 与 $f_X:X_T\to X$ 表示由 $p$ 与 $f$ 作\Term{base-change}{基变换}所得的\Term{morphism}{态射}。
设 $\mathcal F$ 为\Term{quasi-coherent}{拟凝聚}且 $G$-\Term{equivariant}{等变}的 $\mathcal O_X$-\Term{module}{模}，$\theta$ 为 6.5.1（$\star$）的\Term{isomorphism}{同构} $\operatorname{pr}_X^*(\mathcal F)\xrightarrow\sim\lambda^*(\mathcal F)$。由于 $p$ \Term{quasi-compact}{拟紧}且\Term{quasi-separated}{拟分离}，$p_*(\mathcal F)$ \Term{quasi-coherent}{拟凝聚}（参见 EGA I，9.2.1）。

\SourcePage{50}
\begin{SGAresult}{\Term{lemma}{引理}}{6.6.1}
\Term{quasi-coherent}{拟凝聚} $\mathcal O_Y$-\Term{module}{模} $p_*(\mathcal F)$ 是 $G$-\Term{equivariant}{等变}的。
\end{SGAresult}
\begin{proof}
有以下两个\Term{cartesian-square}{笛卡尔方块}：
\[
\begin{tikzcd}
G\times_S X\arrow[r,"{\operatorname{pr}_X}"]\arrow[d,"{q}"']&X\arrow[d,"{p}"]&G\times_S X\arrow[l,"{\lambda}"']\arrow[d,"{q}"]\\
G\times_S Y\arrow[r,"{\operatorname{pr}_Y}"']&Y&G\times_S Y\arrow[l,"{\mu}"].
\end{tikzcd}
\]
因 $p$ \Term{quasi-compact}{拟紧}且\Term{quasi-separated}{拟分离}，$\operatorname{pr}_Y$ 与 $\mu$ \Term{flat}{平坦}，由 EGA III，1.4.15（并结合 EGA $\mathrm{IV}_1$，1.7.21 的补充）可知，$q_*\operatorname{pr}_X^*(\mathcal F)=\operatorname{pr}_Y^*p_*(\mathcal F)$，且 $q_*\lambda^*(\mathcal F)=\mu^*p_*(\mathcal F)$。因此，$\theta$ 诱导\Term{isomorphism}{同构}
\[
\theta_Y:\operatorname{pr}_Y^*p_*(\mathcal F)\xrightarrow\sim\mu^*p_*(\mathcal F),
\]
同样得到 $\theta_Y$ 满足 6.5.1（$\star\star$）的“\Term{cocycle-condition}{余循环条件}”（因为其中所涉及的\Term{morphism}{态射} $G\times_S G\times_S Y\to G\times_S Y$ 都\Term{flat}{平坦}）。这证明了\Term{lemma}{引理}。
\end{proof}
特别地，取 $Y=S$ 并带有 $G$ 的\Term{trivial-action}{平凡作用}。结合注记 6.5.4，便得到 $p_*(\mathcal F)$ 是 $G$-$\mathcal O_S$-\Term{module}{模}。若还有 $G$ 在 $S$ 上\Term{affine}{仿射}，并记 $\mathcal A(G)=\pi_*(\mathcal O_G)$，则由 4.7.2，$p_*(\mathcal F)$ 是 $\mathcal A(G)$-\Term{comodule}{余模}。
另一方面，由 6.4.1，向每个 $f:T\to S$ 赋予
\[
\mathbf W(f_X^*(\mathcal F))(X_T)=\Gamma(X_T,f_X^*(\mathcal F))=\Gamma(T,p_{T*}f_X^*(\mathcal F))
\]
的\Term{functor}{函子} $\prod_{X/S}\mathbf W(\mathcal F)$ 是 $G$-$\mathbf O_S$-\Term{module}{模}。此外，有\Term{canonical}{典范}\Term{morphism}{态射} $\tau:\mathbf W(p_*(\mathcal F))\to\prod_{X/S}\mathbf W(\mathcal F)$；对每个 $f:T\to S$，它由\Term{canonical}{典范}\Term{morphism}{态射}
\[
\Gamma(T,f^*p_*(\mathcal F))\longrightarrow\Gamma(T,p_{T*}f_X^*(\mathcal F))
\]
给出。当\Term{restriction}{限制}到由在 $S$ 上\Term{flat}{平坦}的\Term{scheme}{概形} $T$ 组成的\Term{full-subcategory}{满子范畴}时，它是\Term{isomorphism}{同构}；而且不难核实，$\tau$ 是 $G$-$\mathbf O_S$-\Term{module-morphism}{模态射}。因此得到了以下命题（对于（ii），可比较 [GIT]，第 32 页）。
\begin{SGAresult}{命题}{6.6.2}
设 $S$ 为\Term{scheme}{概形}，$G$ 为\Term{action}{作用}于 $S$-\Term{scheme}{概形} $X$ 的 $S$-\Term{group-scheme}{群概形}，$\mathcal F$ 为\Term{quasi-coherent}{拟凝聚}且 $G$-\Term{equivariant}{等变}的 $\mathcal O_X$-\Term{module}{模}。假设 $\pi:G\to S$ \Term{flat}{平坦}，$p:X\to S$ \Term{quasi-compact}{拟紧}且\Term{quasi-separated}{拟分离}。

（i）则 $p_*(\mathcal F)$ 是\Term{quasi-coherent}{拟凝聚} $G$-$\mathcal O_S$-\Term{module}{模}。另一方面，\Term{canonical}{典范}\Term{morphism}{态射} $\mathbf W(p_*(\mathcal F))\to\prod_{X/S}\mathbf W(\mathcal F)$ 是 $G$-$\mathbf O_S$-\Term{module-morphism}{模态射}；这两个\Term{functor}{函子}在\Term{flat}{平坦} $S$-\Term{scheme}{概形}\Term{category}{范畴}上一致。

（ii）若还有 $G$ 在 $S$ 上\Term{affine}{仿射}，并记 $\mathcal A(G)=\pi_*(\mathcal O_G)$，则 $p_*(\mathcal F)$ 带有 $\mathcal A(G)$-\Term{comodule}{余模}结构。\footnote{编者注（55）：若 $G$ 在 $S$ 上\Term{flat}{平坦}、\Term{quasi-compact}{拟紧}且\Term{quasi-separated}{拟分离}，且 $p_*(\mathcal F)$ 是\Term{flat}{平坦} $\mathcal O_S$-\Term{module}{模}，结论也一样，参见第 $\mathrm{VI_B}$ 讲，11.6.1（ii）。}
\end{SGAresult}

\SourcePage{51}
\textbf{6.7. \Term{stabilizer}{稳定子}。} 设 $S$ 为\Term{scheme}{概形}，$G$ 为\Term{action}{作用}于 $S$-\Term{scheme}{概形} $X$ 的 $S$-\Term{group-scheme}{群概形}，$\mathcal F$ 为\Term{quasi-coherent}{拟凝聚} $G$-\Term{equivariant}{等变} $\mathcal O_X$-\Term{module}{模}。设 $Y$ 为另一个 $S$-\Term{scheme}{概形}，带有 $G$ 的\Term{action}{作用}（可以是\Term{trivial-action}{平凡作用}）；设 $\varphi:Y\to X$ 为一个不必 $G$-\Term{equivariant}{等变}的 $S$-\Term{morphism}{态射}；以 $\operatorname{Stab}_G(\varphi)$ 表示 $\varphi$ 在 $G$ 中的\Term{stabilizer}{稳定子}（参见 6.3.2）。
立即指出（参见第 $\mathrm{VI_B}$ 讲，§6）：若 $X$ 在 $S$ 上\Term{separated}{分离}，$Y$ 在 $S$ 上\Term{essentially-free}{本质自由}（参见该处，定义 6.2.1），则 $\operatorname{Stab}_G(\varphi)$ 可由 $G$ 的一个\Term{closed-subgroup-scheme}{闭子群概形} $H$ 表示。事实上，考虑按\Term{set}{集合}式公式 $r(g,y)=(\varphi(y),g\varphi(g^{-1}y))$ 给出的\Term{morphism}{态射} $r:G\times_S Y\to X\times_S X$，并令 $P=G\times_S Y$，$P'$ 为\Term{diagonal}{对角线} $\Delta_{X/S}$ 在 $r$ 下的\Term{inverse-image}{逆像}。于是有（参见该处，6.2.4（a））
\[
\operatorname{Stab}_G(\varphi)=\prod_{P/G}P'.
\]
所以由该处的结果，若 $X$ 在 $S$ 上\Term{separated}{分离}，$Y$ 在 $S$ 上\Term{essentially-free}{本质自由}，则 $\operatorname{Stab}_G(\varphi)$ 可由 $G$ 的一个\Term{closed-subgroup-scheme}{闭子群概形} $H$ 表示；当 $S$ 是域的\Term{spectrum}{谱}或 $Y=S$ 时，第二个条件自动满足。在这些假设下，$\varphi^*(\mathcal F)$ 是\Term{quasi-coherent}{拟凝聚}且 $H$-\Term{equivariant}{等变}的 $\mathcal O_Y$-\Term{module}{模}（参见 6.5.2）。
因此，若还有 $\pi:G\to S$ 与 $p:Y\to S$ 在 $S$ 上\Term{quasi-compact}{拟紧}且\Term{quasi-separated}{拟分离}，$\pi$ \Term{flat}{平坦}，则由 6.6.2，$p_*\varphi^*(\mathcal F)$ 是 $H$-$\mathcal O_S$-\Term{module}{模}。特别得到：
\begin{SGAresult}{推论}{6.7.1}
设 $S$ 为\Term{scheme}{概形}，$G$ 为\Term{action}{作用}于 $S$-\Term{scheme}{概形} $X$ 的 $S$-\Term{group-scheme}{群概形}，$\mathcal F$ 为\Term{quasi-coherent}{拟凝聚} $G$-\Term{equivariant}{等变} $\mathcal O_X$-\Term{module}{模}。假设 $\pi:G\to S$ \Term{flat}{平坦}，$X\to S$ \Term{quasi-compact}{拟紧}且\Term{separated}{分离}。设 $\tau:S\to X$ 为 $X$ 在 $S$ 上的一个\Term{section}{截面}。

（i）\Term{stabilizer}{稳定子} $H=\operatorname{Stab}_G(\tau)$ 是 $G$ 的\Term{closed-subgroup-scheme}{闭子群概形}，$\tau^*(\mathcal F)$ 为\Term{quasi-coherent}{拟凝聚} $H$-$\mathcal O_S$-\Term{module}{模}。

（ii）若还有 $H$ 在 $S$ 上\Term{affine}{仿射}（例如当 $G$ \Term{affine}{仿射}时），并记 $\mathcal A(H)=\pi_*(\mathcal O_H)$，则 $\tau^*(\mathcal F)$ 是 $\mathcal A(H)$-\Term{comodule}{余模}。
\end{SGAresult}

\textbf{6.8. $G$ 上的 $G$-\Term{equivariant-sheaf}{等变层}。} 最后，指出将在第 II、III 讲中使用的两个结果（以下的 6.8.1 与 6.8.6；特别参见 III，4.25）。
\begin{SGAresult}{命题}{6.8.1}
设 $S$ 为\Term{scheme}{概形}，$\pi:G\to S$ 为 $S$-\Term{group-scheme}{群概形}，$\varepsilon:S\to G$ 为\Term{unit-section}{单位截面}。考虑 $G$ 通过\Term{left-translation}{左平移}在自身上的\Term{action}{作用}。则\Term{functor}{函子} $\mathcal E\mapsto\pi^*(\mathcal E)$ 与 $\mathcal F\mapsto\varepsilon^*(\mathcal F)$ 在\Term{quasi-coherent}{拟凝聚} $\mathcal O_S$-\Term{module}{模}\Term{category}{范畴}与\Term{quasi-coherent}{拟凝聚} $G$-\Term{equivariant}{等变} $\mathcal O_G$-\Term{module}{模}\Term{category}{范畴}之间诱导互为\Term{quasi-inverse}{拟逆}的\Term{equivalence}{等价}。
\end{SGAresult}
\begin{proof}
以 $\mu$ 表示 $G$ 的\Term{multiplication}{乘法}，$\operatorname{pr}_2$ 表示第二\Term{projection}{投影} $G\times_S G\to G$。由于 $\pi\circ\mu=\pi\circ\operatorname{pr}_2$，对每个\Term{quasi-coherent}{拟凝聚} $\mathcal O_S$-\Term{module}{模} $\mathcal E$，有\Term{canonical}{典范}\Term{isomorphism}{同构} $\mu^*\pi^*(\mathcal E)=\operatorname{pr}_2^*\pi^*(\mathcal E)$；容易核实，此\Term{isomorphism}{同构}满足 6.5.1（$\star\star$）的“\Term{cocycle-condition}{余循环条件}”，即 $\pi^*(\mathcal E)$ 为 $G$-\Term{equivariant}{等变} $\mathcal O_G$-\Term{module}{模}。由于 $\varepsilon^*\pi^*(\mathcal E)=\mathcal E$，\Term{functor}{函子} $\mathcal E\mapsto\pi^*(\mathcal E)$
\SourcePage{52}
\Term{fully-faithful}{全忠实}；故只须说明，对每个 $G$-\Term{equivariant}{等变} $\mathcal O_G$-\Term{module}{模} $\mathcal F$，都有 $\mathcal F\simeq\pi^*\varepsilon^*(\mathcal F)$。依假设，有\Term{isomorphism}{同构} $\theta:\operatorname{pr}_2^*(\mathcal F)\xrightarrow\sim\mu^*(\mathcal F)$；沿\Term{set-component}{分量}为 $(\mathrm{id}_G,\varepsilon\circ\pi)$ 的\Term{morphism}{态射} $\tau:G\to G\times_S G$ 取 $\theta$ 的\Term{inverse-image}{逆像}，便得到\Term{isomorphism}{同构} $\mathcal F\xrightarrow\sim\pi^*\varepsilon^*(\mathcal F)$。
\end{proof}
\begin{SGAresult}{注记}{6.8.2}
（a）考虑由 $(g_1,g_2)\cdot g=g_1gg_2^{-1}$ 定义的 $G\times_S G$ 在 $G$ 上的\Term{action}{作用}，则\Term{unit-section}{单位截面} $\varepsilon:S\to G$ 的\Term{stabilizer}{稳定子}是 $G\times_S G$ 的\Term{diagonal-subgroup}{对角子群} $H$。因此，若 $\mathcal F$ 为\Term{quasi-coherent}{拟凝聚} $(G\times_S G)$-\Term{equivariant}{等变} $\mathcal O_G$-\Term{module}{模}，则由 6.5.2，$\mathcal E=\varepsilon^*(\mathcal F)$ 带有 $H$-$\mathcal O_S$-\Term{module}{模}结构。

（b）可以证明，$\mathcal F\mapsto\varepsilon^*(\mathcal F)$ 在\Term{quasi-coherent}{拟凝聚} $(G\times_S G)$-\Term{equivariant}{等变} $\mathcal O_G$-\Term{module}{模}\Term{category}{范畴}与\Term{quasi-coherent}{拟凝聚} $H$-$\mathcal O_S$-\Term{module}{模}\Term{category}{范畴}之间给出\Term{equivalence-of-categories}{范畴等价}。这是更一般的“\Term{descent}{下降}”结果（参见第 IV 讲，§2 及 SGA 1，VIII）的特例，例如参见 [Th87]，1.2--1.3。
\end{SGAresult}
\begin{SGAresult}{注记}{6.8.3}
沿用 6.8.1 的记号。对每个\Term{quasi-coherent}{拟凝聚} $\mathcal O_S$-\Term{module}{模} $\mathcal E$，用 $\pi_*^G\pi^*(\mathcal E)$ 表示 $\pi_*\pi^*(\mathcal E)$ 的子 $\mathcal O_S$-\Term{module}{模}，其在 $S$ 的每个\Term{open-subset}{开集} $V$ 上的\Term{section}{截面}，是满足对每个 $S'\to V$ 与 $g\in G(S')$ 都有 $g\cdot\gamma_{S'}=\gamma_{S'}$ 的 $\gamma\in\Gamma(\pi^{-1}(V),\pi^*(\mathcal E))$。则自然\Term{morphism}{态射} $\mathcal E\to\pi_*^G\pi^*(\mathcal E)$ 是\Term{isomorphism}{同构}：若 $\pi_*\pi^*(\mathcal E)=\mathcal E\otimes_{\mathcal O_S}\pi_*(\mathcal O_G)$（例如当 $G\to S$ \Term{affine}{仿射}，或当 $G\to S$ \Term{quasi-compact}{拟紧}且\Term{quasi-separated}{拟分离}并且 $\mathcal E$ \Term{flat}{平坦}时），这是立即可见的；一般情形也不难核实。
\end{SGAresult}
\begin{SGAresult}{注记}{6.8.4}
设 $S$ 为\Term{scheme}{概形}，$H$ 为\Term{action}{作用}于 $S$-\Term{scheme}{概形} $X$ 的 $S$-\Term{group-scheme}{群概形}，$\mathcal F$ 为 $\mathcal O_X$-\Term{module}{模}。假设 $H$ 在 $S$ 上\Term{flat}{平坦}，用 $\mathbf W_{\mathrm P}(\mathcal F)$ 表示\Term{functor}{函子} $\mathbf W(\mathcal F)$ 在由\Term{flat}{平坦} $S$-\Term{scheme}{概形}组成的\Term{full-subcategory}{满子范畴}上的\Term{restriction}{限制}。由 6.5.1，给 $\mathcal F$ 赋予 $H$-\Term{equivariant}{等变}\Term{module}{模}结构，\Term{equivalence}{等价}于给定 $G\times_S X$ 上的\Term{sheaf}{层}\Term{isomorphism}{同构} $\theta:\operatorname{pr}_X^*(\mathcal F)\to\lambda^*(\mathcal F)$，满足“\Term{cocycle-condition}{余循环条件}”（$\star\star$）；由此可见，要在 $\mathcal F$ 上给定 $H$-\Term{equivariant}{等变}\Term{module}{模}结构，只须在 $\mathbf W_{\mathrm P}(\mathcal F)$ 上给定这样的结构。
\end{SGAresult}
\textbf{回顾 6.8.5。} 设 $X$ 为 $S$-\Term{scheme}{概形}，$Y$ 为 $X$ 的子 $S$-\Term{scheme}{概形}。用 $\mathcal N_{Y/X}$ 表示\Term{immersion}{浸入} $i:Y\hookrightarrow X$ 的\Term{conormal-sheaf}{余法层}（参见 EGA $\mathrm{IV}_4$，16.1.2）。若 $S'\to S$ 为\Term{flat}{平坦}\Term{morphism}{态射}，$i':Y'\hookrightarrow X'$ 表示由 $i$ 作\Term{base-change}{基变换}所得的\Term{immersion}{浸入}，则由该处，16.2.2（iii），有 $\mathcal N_{Y/X}\otimes_{\mathcal O_Y}\mathcal O_{Y'}=\mathcal N_{Y'/X'}$。

\begin{SGAresult}{命题}{6.8.6}
设 $S$ 为\Term{scheme}{概形}，$X$ 为 $S$-\Term{group-scheme}{群概形}，$Y$ 为 $X$ 的子 $S$-\Term{group-scheme}{群概形}。假设 $Y$ 在 $S$ 上\Term{flat}{平坦}。则\Term{conormal-sheaf}{余法层} $\mathcal N_{Y/X}$ 是 $(Y\times_S Y)$-\Term{equivariant}{等变} $\mathcal O_Y$-\Term{module}{模}。
\end{SGAresult}
\begin{proof}
由于 $H=Y\times_S Y$ 在 $S$ 上\Term{flat}{平坦}，由注记 6.8.4，只须在 $\mathbf W_{\mathrm P}(\mathcal N_{Y/X})$ 上赋予 $H$-\Term{equivariant}{等变}\Term{module}{模}结构。设 $S'$ 为\Term{flat}{平坦} $S$-\Term{scheme}{概形}，$Y'\hookrightarrow X'$ 为通过\Term{base-change}{基变换}所得的\Term{immersion}{浸入}，令 $\mathcal N'=\mathcal N_{X/Y}\otimes_{\mathcal O_Y}\mathcal O_{Y'}$。由 6.8.5，有 $\mathcal N'=\mathcal N_{Y'/X'}$。
每个 $h\in H(S')$ 都诱导 $Y'$ 的一个\Term{automorphism}{自同构}，因此对每个 $y\in Y'(S')$，得到\Term{isomorphism}{同构}
\[
\Gamma(S',y^*(\mathcal N'))\xrightarrow\sim\Gamma(S',y^*h^*(\mathcal N')),
\]
\SourcePage{53}
它们在 $\mathbf W_{\mathrm P}(\mathcal N)$ 上给出 $H$-\Term{equivariant}{等变}\Term{module}{模}结构（参见 6.1）。
\end{proof}

\SGAsection{}{参考文献}
\footnote{编者注（56）：本讲所引用的补充参考文献。}
\begin{description}
\item[{[DG70]}] M. Demazure、P. Gabriel，\emph{Groupes Algébriques}（《\Term{algebraic-group}{代数群}》），Masson \& North-Holland，1970。
\item[{[Gr57]}] A. Grothendieck，\emph{Sur quelques points d'algèbre homologique}（《关于\Term{homological-algebra}{同调代数}的若干问题》），Tôhoku Math. J. \textbf{9}（1957），119--221。
\item[{[Ja03]}] J. C. Jantzen，\emph{Representations of algebraic groups}（《\Term{algebraic-group}{代数群}的表示》），Academic Press，1987；第 2 版，Amer. Math. Soc.，2003。
\item[{[GIT]}] D. Mumford，\emph{Geometric invariant theory}（《几何\Term{invariants}{不变量}理论》），Springer-Verlag，1965；第 2 版，与 J. Fogarty 合著，1982；第 3 版，与 J. Fogarty、F. Kirwan 合著，1994。
\item[{[Ni02]}] N. Nitsure，\emph{Representability of $\mathrm{GL}_{\mathcal E}$}（《$\mathrm{GL}_{\mathcal E}$ 的\Term{representability}{可表示性}》），Proc. Indian Acad. Sci. \textbf{112}（2002），第 4 期，539--542。
\item[{[Ni04]}] N. Nitsure，\emph{Representability of Hom implies flatness}（《Hom 的\Term{representability}{可表示性}蕴含\Term{flatness}{平坦性}》），Proc. Indian Acad. Sci. \textbf{114}（2004），第 1 期，7--14。
\item[{[Se68]}] J.-P. Serre，\emph{Groupes de Grothendieck des schémas en groupes réductifs déployés}（《\Term{split-reductive-group}{分裂约化群}\Term{scheme}{概形}的 Grothendieck 群》），Publ. math. I.H.É.S. \textbf{34}（1968），37--52。
\item[{[Th87]}] R. W. Thomason，\emph{Equivariant resolution, linearization, and Hilbert's fourteenth problem over arbitrary base schemes}（《任意\Term{base-scheme}{基概形}上的\Term{equivariant}{等变}\Term{resolution}{消解}、\Term{linear}{线性}化与 Hilbert 第十四问题》），Adv. Maths. \textbf{65}（1987），16--34。
\end{description}
\SourcePage{54}
% 原 PDF 第 54 页为空白页。
